\documentclass[12pt,a4paper]{amsart}
\usepackage{amssymb}
\usepackage[cmtip,all]{xy}
\usepackage{hyperref}

\pagestyle{plain}
\raggedbottom

\textwidth=36pc
\calclayout

\emergencystretch=2em

\newcommand{\+}{\nobreakdash-}
\renewcommand{\:}{\colon}

\newcommand{\rarrow}{\longrightarrow}
\newcommand{\larrow}{\longleftarrow}
\newcommand{\ot}{\otimes}

\newcommand{\lrarrow}{\mskip.5\thinmuskip\relbar\joinrel\relbar
   \joinrel\rightarrow\mskip.5\thinmuskip\relax}
\newcommand{\llarrow}{\mskip.5\thinmuskip\leftarrow\joinrel\relbar
   \joinrel\relbar\mskip.5\thinmuskip}
\newcommand{\bu}{{\text{\smaller\smaller$\scriptstyle\bullet$}}}

\DeclareMathOperator{\Hom}{Hom}
\DeclareMathOperator{\Ext}{Ext}
\DeclareMathOperator{\Spec}{Spec}

\newcommand{\Modl}{{\operatorname{\mathsf{--Mod}}}}
\newcommand{\Qcoh}{{\operatorname{\mathsf{--Qcoh}}}}
\newcommand{\Ctrh}{{\operatorname{\mathsf{--Ctrh}}}}
\newcommand{\Lcth}{{\operatorname{\mathsf{--Lcth}}}}
\newcommand{\Sh}{{\operatorname{\mathsf{--Sh}}}}
\newcommand{\Cosh}{{\operatorname{\mathsf{--Cosh}}}}

\newcommand{\ctra}{{\operatorname{\mathsf{-ctra}}}}

\newcommand{\Ab}{\mathsf{Ab}}

\newcommand{\sop}{\mathsf{op}}

\newcommand{\prj}{\mathsf{prj}}
\newcommand{\inj}{\mathsf{inj}}
\newcommand{\fl}{\mathsf{fl}}
\renewcommand{\cot}{\mathsf{cot}}
\newcommand{\vfl}{\mathsf{vfl}}
\newcommand{\cta}{\mathsf{cta}}
\newcommand{\qs}{\mathsf{qs}}
\newcommand{\lct}{\mathsf{lct}}
\newcommand{\lin}{\mathsf{lin}}

\newcommand{\bB}{\mathbf B}
\newcommand{\bW}{\mathbf W}
\newcommand{\bT}{\mathbf T}

\newcommand{\cO}{\mathcal O}
\newcommand{\cF}{\mathcal F}
\newcommand{\cI}{\mathcal I}
\newcommand{\cM}{\mathcal M}
\newcommand{\cN}{\mathcal N}

\newcommand{\fF}{\mathfrak F}
\newcommand{\fG}{\mathfrak G}
\newcommand{\fI}{\mathfrak I}
\newcommand{\fJ}{\mathfrak J}
\newcommand{\fK}{\mathfrak K}
\newcommand{\fL}{\mathfrak L}
\newcommand{\fP}{\mathfrak P}
\newcommand{\fQ}{\mathfrak Q}
\newcommand{\fR}{\mathfrak R}

\newcommand{\m}{\mathfrak m}
\newcommand{\p}{\mathfrak p}

\newcommand{\sK}{\mathsf K}
\newcommand{\sE}{\mathsf E}
\newcommand{\sF}{\mathsf F}
\newcommand{\sC}{\mathsf C}
\newcommand{\sS}{\mathsf S}

\newcommand{\boQ}{\mathbb Q}
\newcommand{\boZ}{\mathbb Z}


\newcommand{\Section}[1]{\bigskip\section{#1}\medskip}
\setcounter{tocdepth}{1}

\theoremstyle{plain}
\newtheorem{thm}{Theorem}[section]
\newtheorem{prop}[thm]{Proposition}
\newtheorem{lem}[thm]{Lemma}
\newtheorem{cor}[thm]{Corollary}
\theoremstyle{definition}
\newtheorem{rem}[thm]{Remark}
\newtheorem{rems}[thm]{Remarks}
\newtheorem{quest}[thm]{Question}
\newtheorem{ex}[thm]{Example}
\newtheorem{exs}[thm]{Examples}

\begin{document}

\title{The contraherent version of the theorem \\
of Sl\'avik and \v S\v tov\'\i\v cek}

\author{Leonid Positselski}

\address{Institute of Mathematics, Czech Academy of Sciences \\
\v Zitn\'a~25, 115~67 Prague~1 \\ Czech Republic} 

\email{positselski@math.cas.cz}

\begin{abstract}
 Let $X$ be a quasi-compact, quasi-separated scheme that is \emph{not}
semi-separated.
 We present a locally cotorsion contraherent cosheaf on $X$ that
\emph{does not} have an admissible monomorphism into any locally
injective locally contraherent cosheaf.
 The construction and proof follow the arguments of Sl\'avik
and \v St\!'ov\'\i\v cek.
 A complementary positive result, claiming that every $\bW$\+locally
contraherent cosheaf on a Noetherian scheme $X$ of finite Krull
dimension has an admissible monomorphism into a locally cotorsion
$\bW$\+locally contraherent cosheaf, can be found in
the preprint~\cite[Section~4.7]{Pform}.
\end{abstract}

\maketitle

\tableofcontents


\section{Introduction}
\medskip

 Various classes of adjusted objects, such the classical flat,
projective, and injective modules, are used in homological algebra
in order to build resolutions and construct derived functors.
 There are always enough projective and injective objects in
the category of modules over an associative ring.
 The situation becomes more complicated when one passes to
algebraic geometry and considers the category of quasi-coherent
sheaves over a scheme.

 For any scheme $X$, the category of quasi-coherent sheaves $X\Qcoh$
is a Grothendieck abelian category~\cite[Section Tag~077K]{SP};
so it has all limits and enough injective objects.
 However, in any abelian category with infinite products and enough
projective objects, the direct product functors are exact.
 The infinite products of quasi-coherent sheaves are usually \emph{not}
exact functors~\cite[Example~4.9]{Kra}, \cite{Kan}; hence the category
$X\Qcoh$ cannot have enough projective objects (and in fact, usually
has no nonzero projective objects at all).  {\hbadness=1600\par}

 The conventional approach is to use flat quasi-coherent sheaves in
lieu of the nonexistent projective ones~\cite{M-th}.
 On a quasi-compact, semi-separated scheme $X$, any quasi-coherent sheaf
is a quotient sheaf of a flat quasi-coherent
sheaf~\cite[Section~2.4]{M-n}, \cite[Section~3.2]{M-th}.
 Another and perhaps clearer proof of this result can be found
in~\cite[Lemma~A.1]{EP}.
 Essentially the same argument as in~\cite{EP} can be used to prove
a stronger assertion that any quasi-coherent sheaf on $X$ is
a quotient sheaf of a \emph{very flat} quasi-coherent
sheaf~\cite[Lemma~4.1.1]{Pcosh}.

 On the other hand, let $X$ be a quasi-compact, quasi-separated scheme
that is \emph{not} semi-separated.
 Then the construction of Sl\'avik and
\v St\!'ov\'\i\v cek~\cite[Theorem~2.2]{SS} presents an example of
a quasi-coherent sheaf on $X$ that is \emph{not} a quotient sheaf of
any flat quasi-coherent sheaf.

 The \emph{contraherent cosheaves} are the dual analogues of
quasi-coherent sheaves~\cite[Preface]{Pcosh}, \cite{Pphil}.
 For any scheme $X$, there is the exact (but not abelian!) category
of contraherent cosheaves $X\Ctrh$.
 Choosing an open covering $\bW$ for the scheme $X$, one also obtains
a larger exact category of \emph{$\bW$\+locally contraherent cosheaves}
$X\Lcth_\bW$.
 The categories $X\Ctrh$ and $X\Lcth_\bW$ have exact functors of
infinite products and, under mild assumptions on the scheme $X$,
enough projective objects~\cite[Section~0.18, Lemma~4.5.1,
and Corollary~6.2.4]{Pcosh}.

 Dual-analogously to the categories of quasi-coherent sheaves not
having enough projective objects, one does not expect the categories
of contraherent or locally contraherent cosheaves to have enough
injective objects (though we are not aware of any specific
counterexamples).
 The \emph{locally injective} contraherent or locally contraherent
cosheaves are the closest dual analogues of flat (or very flat)
quasi-coherent sheaves.

 On a quasi-compact semi-separated scheme $X$, any $\bW$\+locally
contraherent cosheaf is an admissible subobject of a locally injective
$\bW$\+locally contraherent cosheaf~\cite[Lemma~4.3.2(a)]{Pcosh}.
 The proof of this result is dual-analogous to the proofs of
the existence of enough flat or very flat quasi-coherent sheaves
given in~\cite[Lemma~A.1]{EP} and~\cite[Lemma~4.1.1]{Pcosh}.

 Now let $X$ be again a quasi-compact, quasi-separated scheme that is
not semi-separated.
 Then very close dual analogues of the construction and arguments of
Sl\'avik and \v St\!'ov\'\i\v cek~\cite[Section~2]{SS} provide
an example of a contraherent cosheaf $\fG$ on $X$ that is \emph{not}
an admissible subobject of any locally injective locally contraherent
cosheaf.
 This is the main (negative) result of the present paper.

 The (locally) contraherent cosheaves forming the exact categories
$X\Ctrh$ and $X\Lcth_\bW$ are glued from what we call
\emph{contraadjusted modules} over commutative
rings~\cite[Section~1.1]{Pcosh}, \cite{ST}, \cite{Pcta}; so such
cosheaves are called \emph{locally contraadjusted}.
 This is the widest reasonable class of modules to glue cosheaves from.
 The \emph{locally cotorsion} contraherent or locally contraherent
cosheaves, glued from \emph{cotorsion modules}~\cite{En2},
form intermediate exact categories between the exact categories of
locally contraadjusted and locally injective (locally) contraherent
cosheaves.

 Our contraherent cosheaf $\fG$ mentioned two paragraphs above is,
actually, a locally cotorsion contraherent cosheaf.
 So embedding locally cotorsion contraherent cosheaves into locally
injective ones appears to be no easier task than embedding locally
contraadjusted contraherent cosheaves into locally injective ones.

 On the positive side, let $X$ be a Noetherian scheme of finite Krull
dimension.
 Notice that such a scheme is always quasi-compact and quasi-separated,
but \emph{need not} be semi-separated.
 Then, according to~\cite[Proposition~4.7.11]{Pform}, any (i.~e., any
locally contraadjusted) $\bW$\+locally contraherent cosheaf on $X$ is
an admissible subobject of a locally cotorsion $\bW$\+locally
contraherent cosheaf.

\subsection*{Acknowledgement}
 The author is supported by the GA\v CR project 26-22734S
and the Institute of Mathematics, Czech Academy of Sciences
(research plan RVO:~67985840).


\Section{Category-Theoretic Preliminaries}

 We suggest the survey paper~\cite{Bueh} as the background reference
source on exact categories (in the sense of Quillen).
 The assumption of \emph{weak
idempotent-completeness}~\cite[Section~7]{Bueh} simplifies
the theory of exact categories.
 For the purposes of the present paper, one can safely assume all
the exact categories to satisfy the even stronger
\emph{idempotent-completeness} condition~\cite[Section~6]{Bueh}.

 Given an exact category $\sK$ and a full additive subcategory
$\sE\subset\sK$ closed under extensions in $\sK$, the category $\sE$
is endowed with the \emph{inherited} (or \emph{induced}) exact
category structure.
 This means that the admissible short exact sequences in $\sE$ are
the admissible short exact sequences in $\sK$ with the terms belonging
to~$\sE$.

 The Yoneda Ext groups $\Ext_\sE^n(X,Y)$, \,$X$, $Y\in\sE$,
\,$n\in\boZ_{\ge0}$ in an exact category $\sE$ can be defined in
the way similar to the classical construction for abelian categories.
 In particular, $\Ext_\sE^1(X,Y)$ is the group of equivalence
classes of admissible short exact sequences $0\rarrow Y\rarrow Z
\rarrow X\rarrow0$ in~$\sE$.

 A discussion of projective and injective objects in exact categories
can be found in~\cite[Section~11]{Bueh}.
 We denote the full subcategory of projective objects by
$\sE_\prj\subset\sE$ and the full subcategory of injective objects
by $\sE^\inj\subset\sE$.

\subsection{Resolving and coresolving subcategories}
\label{resolving-subsecn}
 Let $\sE$ be an exact category.
 A full subcategory $\sF\subset\sE$ is said to be \emph{generating}
in $\sE$ if for every object $E\in\sE$ there exists an object
$F\in\sF$ together with an admissible epimorphism $F\rarrow E$.
 Dually, a full subcategory $\sC\subset\sE$ is said to be
\emph{cogenerating} in $\sE$ if for every object $E\in\sE$ there
exists an object $C\in\sC$ together with an admissible monomorphism
$E\rarrow C$.

 A generating full subcategory $\sF\subset\sE$ is said to be
\emph{resolving} if $\sF$ is closed under extensions and kernels of
admissible epimorphisms in~$\sE$.
 Dually, a cogenerating full subcategory $\sC\subset\sE$ is said to be
\emph{coresolving} if $\sC$ is closed under extensions and cokernels
of admissible monomorphisms in~$\sE$.

 For example, the full subcategory of projective objects $\sE_\prj$
is resolving in $\sE$ if and only if it is generating in $\sE$, and
if and only if there are enough projectives in~$\sE$.
 Dually, the full subcategory of injective objects $\sE^\inj$ is
coresolving in $\sE$ if and only if it is cogenerating in $\sE$, and
if and and only if there are enough injectives in~$\sE$.
 To give yet another example, for any associative ring $R$,
the full subcategory of flat $R$\+modules $R\Modl_\fl$ is resolving
in the abelian category of left $R$\+modules $R\Modl$.

 Let $\sE$ be a weakly idempotent-complete exact category,
$\sF\subset\sE$ be a resolving subcategory, $E\in\sE$ be an object,
and $d\ge-1$ be an integer.
 One says that the \emph{$\sF$\+resolution dimension} of $E$ does not
exceed~$d$ if there exists an exact sequence $0\rarrow F_d\rarrow
F_{d-1}\rarrow\dotsb\rarrow F_0\rarrow E\rarrow0$ in $\sE$ with
the objects $F_i\in\sF$ for all $0\le i\le d$.
 Dually, given a coresolving subcategory $\sC\subset\sE$, one says that
the \emph{$\sC$\+coresolution dimension} of an object $E\in\sE$ does
not exceed~$d$ if there exists an exact sequence $0\rarrow E\rarrow
C^0\rarrow C^1\rarrow\dotsb\rarrow C^d\rarrow0$ in $\sE$ with
the objects $C^i\in\sC$ for all $0\le i\le d$.

 It is a classical observation in homological algebra that
the projective/injective dimension of an object does not depend on
the choice of a projective/injective (co)resolution.
 Similarly, the flat dimension of an $R$\+module does not depend on
the choice of a flat resolution.
 Such results, stated in the following proposition, are actually valid
for all (co)resolving subcategories in (weakly idempotent-complete)
exact categories.

\begin{prop} \label{co-resolution-dimension-prop}
\textup{(a)} Let\/ $\sE$ be a weakly idempotent-complete exact category
and\/ $\sF\subset\sE$ be a resolving subcategory.
 Let $E\in\sE$ be an object of\/ $\sF$\+resolution dimension\/~$\le d$
and\/ $0\rarrow F\rarrow F_{d-1}\rarrow\dotsb\rarrow F_0\rarrow E
\rarrow0$ be an exact sequence in\/ $\sE$ with the objects $F_i\in\sF$
for all\/ $0\le i\le d-1$.
 Then one has $F\in\sF$. \par
\textup{(b)} Let\/ $\sE$ be a weakly idempotent-complete exact category
and\/ $\sC\subset\sE$ be a coresolving subcategory.
 Let $E\in\sE$ be an object of\/ $\sC$\+coresolution dimension\/~$\le d$
and\/ $0\rarrow E\rarrow C^0\rarrow\dotsb\rarrow C^{d-1}\rarrow C
\rarrow0 $ be an exact sequence in\/ $\sE$ with the objects $C_i\in\sC$
for all\/ $0\le i\le d-1$.
 Then one has $C\in\sC$.
\end{prop}

\begin{proof}
 Various expositions on various generality levels can be found
in~\cite[Lemma~2.1]{Zhu}, \cite[Proposition~2.3(1)]{Sto},
and~\cite[Corollary~A.5.2]{Pcosh}.  \hbadness=1600
\end{proof}

\subsection{Cotorsion pairs}
 The notion of a \emph{cotorsion pair} (also known as a \emph{cotorsion
theory}) goes back to Salce~\cite{Sal}.
 The main result was obtained in the paper of Eklof and
Trlifaj~\cite{ET}.
 The book~\cite[Chapter~6]{GT} is the main reference source in
the module-theoretic context, while the papers~\cite[Sections~8
and~12]{Pcta} and~\cite[Introduction]{Pctrl} can be used for
introductory reading.
 Various category-theoretic expositions can be found in
the papers~\cite{SaoSt,Pal}, the lecture notes~\cite{Sto-ICRA},
the announcement~\cite{Pphil}, and the book
manuscript~\cite[Appendix~B]{Pcosh}.

 Let $\sE$ be an exact category and $\sF$, $\sC\subset\sE$ be two
classes of objects.
 One denotes by $\sF^{\perp_1}\subset\sE$ the class of all objects
$X\in\sE$ such that $\Ext^1_\sE(F,X)=0$ for all $F\in\sF$.
 Dually, the notation ${}^{\perp_1}\sC\subset\sE$ stands for the class
of all objects $Y\in\sE$ such that $\Ext^1_\sE(Y,C)=0$ for all
$C\in\sC$.

 A pair of classes of objects $(\sF,\sC)$ in $\sE$ is called
a \emph{cotorsion pair} if one has $\sC=\sF^{\perp_1}$ and
$\sF={}^{\perp_1}\sC$.
 For any class of objects $\sS\subset\sE$, the pair of classes
$\sC=\sS^{\perp_1}$ and $\sF={}^{\perp_1}\sC$ is a cotorsion pair
in~$\sE$.
 The latter cotorsion pair is said to be \emph{generated by}
the class of objects $\sS\subset\sE$.

 For example, for any exact category $\sE$, the pair of classes of
objects $(\sE_\prj,\sE)$ is a cotorsion pair in $\sE$, called
the \emph{projective cotorsion pair}.
 Dually, the pair of classes of objects $(\sE,\sE^\inj)$ is also
a cotorsion pair in $\sE$, called the \emph{injective cotorsion pair}.

 Notice that there is \emph{no} assumption of existence of enough
projective or injective objects in $\sE$ in the previous paragraph.
 So what does it mean, in application to a cotorsion pair, that
there are enough objects in the classes $\sF$ and~$\sC$\,?
 The rather nontrivial answer to this question is provided by
the following definition of a \emph{complete} cotorsion pair in
an exact category.

 A cotorsion pair $(\sF,\sC)$ in $\sE$ is said to be complete if,
for every object $E\in\sE$, there exist (admissible) short exact
sequences
\begin{alignat}{4}
 0&\lrarrow C'&&\lrarrow F&&\lrarrow E&&\lrarrow0,
\label{special-precover-sequence} \\
 0&\lrarrow E&&\lrarrow C&&\lrarrow F'&&\lrarrow0
\label{special-preenvelope-sequence}
\end{alignat}
with some objects $F$, $F'\in\sF$ and $C$, $C'\in\sC$.

 The short exact sequence~\eqref{special-precover-sequence} is called
a \emph{special precover sequence}.
 The short exact sequence~\eqref{special-preenvelope-sequence} is
called a \emph{special preenvelope sequence}.
 Collectively, the sequences~(\ref{special-precover-sequence}\+-%
\ref{special-preenvelope-sequence}) are referred to as
the \emph{approximation sequences}.


\Section{Module-Theoretic Preliminaries}

 Given an associative ring $R$, we denote by $R\Modl$ the abelian
category of left $R$\+modules.
 Accordingly, $R\Modl_\prj$ and $R\Modl^\inj\subset R\Modl$ are the full
subcategories of projective and injective $R$\+modules, respectively.
 The notation $R\Modl_\fl\subset R\Modl$ for the full subcategory
of flat $R$\+modules was already mentioned in
Section~\ref{resolving-subsecn}.

\subsection{Cotorsion modules} \label{cotorsion-modules-subsecn}
 A left $R$\+module $F$ is said to be \emph{cotorsion} (in the sense
of Enochs~\cite{En2}) if one has $\Ext_R^1(F,C)=0$ for all flat
left $R$\+modules~$F$.
 This is equivalent to having $\Ext_R^n(F,C)=0$ for all flat left
$R$\+modules $F$ and all integers $n\ge1$ (as one can easily see
using the facts that projective modules are flat and the kernels of
subjective morphisms of flat modules are flat).

 We denote the full subcategory of cotorsion $R$\+modules by
$R\Modl^\cot\subset R\Modl$.
 The full subcategory $R\Modl^\cot$ is closed under extensions,
cokernels of monomorphisms, and infinite products in $R\Modl$,
while the full subcategory $R\Modl_\fl$ is well known to be closed
under extensions, kernels of epimorphisms, and direct limits in
$R\Modl$.
 So both the full subcategories $R\Modl^\cot$ and $R\Modl_\fl$
inherit exact category structures from the abelian exact structure
of $R\Modl$.

\begin{thm} \label{flat-cotorsion-pair}
 For any associative ring $R$, the pair of classes (flat $R$\+modules,
cotorsion $R$\+modules) is a complete cotorsion pair in $R\Modl$.
 In other words: \par
\textup{(a)} for any $R$\+module $M$ there exists a short exact
sequence of $R$\+modules\/ $0\rarrow C\rarrow F\rarrow M\rarrow0$
with a flat $R$\+module $F$ and a cotorsion $R$\+module~$C$; \par
\textup{(b)} for any $R$\+module $M$ there exists a short exact
sequence of $R$\+modules\/ $0\rarrow M\rarrow C\rarrow F\rarrow0$
with a cotorsion $R$\+module $C$ and a flat $R$\+module~$F$.
\end{thm}

\begin{proof}
 The assertion that $(R\Modl_\fl,\allowbreak\>R\Modl^\cot)$ is
\emph{a cotorsion pair} already requires a proof.
 One has $R\Modl^\cot=(R\Modl_\fl)^{\perp_1}$ by the definition; but
why does the equality $R\Modl_\fl={}^{\perp_1}(R\Modl^\cot)$ hold?
 A proof of that can be found in~\cite[Lemma~3.4.1]{Xu}.
 \emph{Completeness} of this cotorsion pair is harder to prove;
the argument, based on~\cite[Theorems~2 and~10]{ET}, appeared
in the paper~\cite[Section~2]{BBE}.
\end{proof}

 The cotorsion pair $(R\Modl_\fl,\allowbreak\>R\Modl^\cot)$ is called
the \emph{flat cotorsion pair} in $R\Modl$.

\begin{lem} \label{flat-cotorsion-injective-tensor-Hom-lemma}
 Let $R$ be a commutative ring. \par
\textup{(a)} For any two flat $R$\+modules $F$ and $G$,
the $R$\+module $F\ot_RG$ is flat. \par
\textup{(b)} For any flat $R$\+module $F$ and any cotorsion
$R$\+module $C$, the $R$\+module\/ $\Hom_R(F,C)$ is cotorsion. \par
\textup{(c)} For any $R$\+module $M$ and any injective
$R$\+module $J$, the $R$\+module\/ $\Hom_R(M,J)$ is cotorsion. \par
\textup{(d)} For any flat $R$\+module $F$ and any injective
$R$\+module $J$, the $R$\+module\/ $\Hom_R(F,J)$ is injective.
\end{lem}

\begin{proof}
 All the assertions are well known.
 Part~(a) is easy.
 Part~(c) can be found in~\cite[Lemma~2.1]{En2}.
 Parts~(b\+-d) can be found in~\cite[Lemma~1.3.2]{Pcosh};
see also the more general~\cite[Lemma~1.5]{Pdomc} with its
slightly different proof.
\end{proof}

 The next lemma is even more well-known and easy.

\begin{lem} \label{injective-change-of-scalars-lemma}
 Let $R\rarrow S$ be a homomorphism of commutative rings. \par
\textup{(a)} If $S$ is a flat $R$\+module, then any injective
$S$\+module is injective as an $R$\+module. \par
\textup{(b)} For any injective $R$\+module $J$, the $S$\+module\/
$\Hom_R(S,J)$ is injective. \qed
\end{lem}

\begin{lem} \label{cotorsion-change-of-scalars-lemma}
 Let $R\rarrow S$ be a homomorphism of commutative rings. \par
\textup{(a)} Any cotorsion $S$\+module is cotorsion as an $R$\+module.
\par
\textup{(b)} If $S$ is a flat $R$\+module, then for any cotosion
$R$\+module $C$, the $S$\+module\/ $\Hom_R(S,C)$ is cotorsion.
\end{lem}

\begin{proof}
 This is a particular case of~\cite[Lemmas~1.3.4(a) and~1.3.5(a)]{Pcosh}
or~\cite[Lemma~1.6(a,c)]{Pdomc}.
\end{proof}

\subsection{Contraadjusted modules}
 Now let $R$ be a commutative ring.
 For any element $s\in R$, we denote by $R[s^{-1}]$ the localization
of the ring $R$ at the element~$s$, i.~e., the ring obtained from $R$
by formally inverting~$s$.
 So, in other words, $R[s^{-1}]=S^{-1}R$, where $S$ is
the multiplicative subset $S=\{1,s,s^2,s^3,\dotsc\}\subset R$.

 An $R$\+module $C$ is said to be
\emph{contraadjusted}~\cite[Section~1.1]{Pcosh}, \cite[Sections~2
and~5]{ST}, \cite[Sections~2 and~8]{Pcta}, \cite[Section~4.3]{Pphil},
\cite[Section~2]{Pal} if $\Ext^1_R(R[s^{-1}],C)=0$ for all
elements $s\in R$.
 It is helpful to keep in mind that the projective dimension of
the $R$\+module $R[s^{-1}]$ never exceeds~$1$ \,\cite[proof of
Lemma~2.2]{ST}, \cite[proof of Lemma~2.1]{Pcta}; that is why no
higher $\Ext$ vanishing condition is needed in this definition.
 Any quotient $R$\+module of a contraadjusted $R$\+module is
contraadjusted.

 An $R$\+module $F$ is said to be \emph{very
flat}~\cite[Section~1.1]{Pcosh}, \cite[Section~2]{ST},
\cite[Section~0.5]{PSl1}, \cite[Section~4.3]{Pphil}
if $\Ext^1_R(F,C)=0$ for all contraadjusted $R$\+modules~$C$.
 The projective dimensions of very flat $R$\+modules do not
exceed~$1$ \,\cite[Section~1.1]{Pcosh}, \cite[Lemma~2.2]{ST},
\cite[Theorem~4.5(d)]{Pphil}.

\begin{ex} \label{open-affine-very-flat-example}
 Let $U$ be an affine scheme and $V\subset U$ be an affine open
subscheme.
 Then the $\cO(U)$\+module $\cO(V)$ is very
flat~\cite[Lemma~1.2.4]{Pcosh}.

 This is a special case of a much more general (and much more
difficult) result of~\cite[Main Theorem~1.1]{PSl1}
(see also a discussion in~\cite[Remark~5.1]{Pphil}).
\end{ex}

 We denote the full subcategory of contraadjusted $R$\+modules by
$R\Modl^\cta\subset R\Modl$ and the full subcategory of very flat
$R$\+modules by $R\Modl_\vfl\subset R\Modl$.
 The full subcategory $R\Modl^\cta$ is closed under extensions,
quotients, and infinite products, while the full subcategory
$R\Modl_\vfl$ is closed under extensions, kernels of epimorphisms,
and infinite direct sums in $R\Modl$.
 So both the full subcategories $R\Modl^\cta$ and $R\Modl_\vfl$
inherit exact category structures from the abelian exact structure
of $R\Modl$.

\begin{thm} \label{very-flat-cotorsion-pair}
 For any commutative ring $R$, the pair of classes (very flat
$R$\+modules, contraadjusted $R$\+modules) is a complete cotorsion
pair in $R\Modl$.
 In other words: \par
\textup{(a)} for any $R$\+module $M$ there exists a short exact
sequence of $R$\+modules\/ $0\rarrow C\rarrow F\rarrow M\rarrow0$ with
a very flat $R$\+module $F$ and a contraadjusted $R$\+module~$C$; \par
\textup{(b)} for any $R$\+module $M$ there exists a short exact
sequence of $R$\+modules\/ $0\rarrow M\rarrow C\rarrow F\rarrow0$ with
a contraadjusted $R$\+module $C$ and a very flat $R$\+module~$F$.
\end{thm}

\begin{proof}
 This is a special case of the Eklof--Trlifaj
theorem~\cite[Theorems~2 and~10]{ET}, \cite[Theorem~6.11]{GT}.
 Some details specific to the situation at hand can be found
in~\cite[Theorem~1.1.1]{Pcosh}, \cite{ST}, \cite[Theorem~4.5]{Pphil}.
\end{proof}

 The word ``contraadjusted'' means ``adjusted to contraherent
cosheaves''.
 The cotorsion pair $(R\Modl_\vfl,\allowbreak\>R\Modl^\cta)$ is called
the \emph{very flat cotorsion pair} in $R\Modl$.

 The following lemma is the very flat/contraadjusted version of
Lemma~\ref{flat-cotorsion-injective-tensor-Hom-lemma}.

\begin{lem} \label{vfl-cta-tensor-Hom-lemma}
 Let $R$ be a commutative ring. \par
\textup{(a)} For any two very flat $R$\+modules $F$ and $G$,
the $R$\+module $F\ot_RG$ is very flat. \par
\textup{(b)} For any very flat $R$\+module $F$ and any contraadjusted
$R$\+module $C$, the $R$\+module\/ $\Hom_R(F,C)$ is contraadjusted.
\end{lem}

\begin{proof}
 This is~\cite[Lemma~1.2.1]{Pcosh}.
\end{proof}

\begin{lem} \label{very-flat-change-of-scalars-lemma}
 Let $R\rarrow S$ be a homomorphism of commutative rings. \par
\textup{(a)} If the $R$\+module $S[s^{-1}]$ is very flat for every
element $s\in S$, then any very flat $S$\+module is very flat as
an $R$\+module. \par
\textup{(b)} For any very flat $R$\+module $F$, the $S$\+module
$S\ot_RF$ is very flat.
\end{lem}

\begin{proof}
 This is~\cite[Lemmas~1.2.3(b) and~1.2.2(b)]{Pcosh}.
\end{proof}

 The next lemma is the contraadjusted version of
Lemma~\ref{cotorsion-change-of-scalars-lemma}.

\begin{lem} \label{contraadjusted-change-of-scalars-lemma}
 Let $R\rarrow S$ be a homomorphism of commutative rings. \par
\textup{(a)} Any contraadjusted $S$\+module is also contraadjusted
as an $R$\+module. \par
\textup{(b)} If the $R$\+module $S[s^{-1}]$ is very flat for every
element $s\in S$, then the $S$\+module\/ $\Hom_R(S,C)$ is
contraadjusted for any contraadjusted $R$\+module~$C$.
\end{lem}

\begin{proof}
 This is~\cite[Lemmas~1.2.2(a) and~1.2.3(a)]{Pcosh}.
\end{proof}

\begin{lem} \label{Hom-into-pushout-lemma}
 Let $R$ be an associative ring and
\begin{equation} \label{pushout-diagram-of-modules}
\begin{gathered}
 \xymatrix{
  0 \ar[r] & K \ar[r] \ar[d] & L \ar[r] \ar[d] & M \ar[r] & 0 \\
  0 \ar[r] & K' \ar[r] & L' \ar[ur]
 }
\end{gathered}
\end{equation}
be a pushout diagram of short exact sequences of left $R$\+modules
(so the upper line is a short exact sequence and the square on
the left-hand side is a pushout square).
 In this context: \par
\textup{(a)} If $K$ is a cotorsion $R$\+module, and $F$ is a flat
left $R$\+module, then applying the functor\/ $\Hom_R(F,{-})$ to
the diagram~\eqref{pushout-diagram-of-modules} produces
a pushout diagram of abelian groups. \par
\textup{(b)} If the ring $R$ is commutative, $K$ is a contraadjusted
$R$\+module, and $F$ is a very flat $R$\+module, then applying
the functor\/ $\Hom_R(F,{-})$ to
the diagram~\eqref{pushout-diagram-of-modules} produces
a pushout diagram of $R$\+modules.
\end{lem}

\begin{proof}
 In both parts~(a) and~(b), the point is that the short exact
sequence $0\rarrow K\rarrow L\oplus K'\rarrow L'\rarrow0$ remains
exact after the functor $\Hom_R(F,{-})$ is applied, due to
the adjustedness assumptions imposed on $K$ and~$F$.
 See, e.~g., \cite[Proposition~2.12]{Bueh} for a background
discussion.
\end{proof}

\subsection{Locality lemmas}
 Throughout this section, we consider a commutative ring $R$ and
a finite collection of homomorphisms of commutative rings
$R\rarrow S_\alpha$, \,$1\le\alpha\le N$, such that the related
finite collection of morphisms of affine schemes $\Spec S_\alpha
\rarrow\Spec R$ is an open covering of $\Spec R$.
 So, in particular, for every index~$\alpha$, the morphism
$\Spec S_\alpha\rarrow\Spec R$ is an open immersion of schemes.

\begin{lem} \label{very-flatness-local}
\textup{(a)} An $R$\+module $F$ is flat if and only if
the $S_\alpha$\+module $S_\alpha\ot_RF$ is flat for every
index\/~$\alpha$. \par
\textup{(b)} An $R$\+module $F$ is very flat if and only if
the $S_\alpha$\+module $S_\alpha\ot_RF$ is very flat for every
index\/~$\alpha$.
\end{lem}

\begin{proof}
 Part~(a) is well known, while part~(b) is~\cite[Lemma~1.2.6(a)]{Pcosh}
or~\cite[Example~2.5]{Pal}.
 The ``only if'' assertion of part~(b) holds by
Lemma~\ref{very-flat-change-of-scalars-lemma}(b).
 The ``if'' implications of both parts~(a) and~(b) are provable using
the finite \v Cech coresolution
\begin{multline} \label{module-cech-coresolution}
 0\lrarrow F\lrarrow\bigoplus\nolimits_{\alpha=1}^N
 S_\alpha\ot_RF\lrarrow\bigoplus\nolimits_{1\le\alpha<\beta\le N}
 S_\alpha\ot_RS_\beta\ot_RF \\
 \lrarrow\dotsb\lrarrow S_1\ot_R\dotsb\ot_RS_N\ot_RF\lrarrow0,
\end{multline}
which is an exact sequence of $R$\+modules for any $R$\+module~$F$.
 To prove the ``if'' implication in~(b), one needs to
use~\eqref{module-cech-coresolution} together with
Example~\ref{open-affine-very-flat-example},
Lemma~\ref{very-flat-change-of-scalars-lemma}(a\+-b), and
the fact that the class of very flat $R$\+modules is closed
under kernels of epimorphisms.
 One proceeds by induction, moving from the rightmost to
the leftmost end of the sequence~\eqref{module-cech-coresolution}
and proving that its $R$\+modules of cocycles are very flat.
\end{proof}

\begin{lem} \label{contraadjusted-cech-resolution}
 Let $C$ be a contraadjusted $R$\+module.
 Then the finite \v Cech resolution
\begin{multline} \label{cta-module-cech-resolution}
 0\lrarrow\Hom_R(S_1\ot_R\dotsb\ot_R S_N,\>C)\lrarrow\dotsb \\
 \lrarrow\bigoplus\nolimits_{1\le\alpha<\beta\le N}
 \Hom_R(S_\alpha\ot_RS_\beta,\>C)\lrarrow
 \bigoplus\nolimits_{\alpha=1}^N\Hom_R(S_\alpha,C)
 \lrarrow C\lrarrow 0
\end{multline}
is an exact sequence of $R$\+modules (and in fact, an exact
sequence in the exact category of contraadjusted $R$\+modules
$R\Modl^\cta$).
\end{lem}

\begin{proof}
 This is~\cite[Lemma~1.2.6(b)]{Pcosh}.
 Set $F=R$; then the sequence~\eqref{module-cech-coresolution}
is exact in the exact category $R\Modl_\vfl$ (as per the proof of
Lemma~\ref{very-flatness-local}) above.
 Hence, applying the functor $\Hom_R({-},C)$
to~\eqref{module-cech-coresolution} produces an exact sequence
in $R\Modl^\cta$ in view of Lemma~\ref{vfl-cta-tensor-Hom-lemma}(b).
 This is the desired exact sequence~\eqref{cta-module-cech-resolution}.
\end{proof}

 The following lemma is a dual-analogous version of
Lemma~\ref{very-flatness-local}.

\begin{lem} \label{cotorsion-injective-colocal-if-contraadjusted}
 Let $C$ be a contraadjusted $R$\+module. \par
\textup{(a)} The $R$\+module $C$ is cotorsion if and only if
the $S_\alpha$\+module\/ $\Hom_R(S_\alpha,C)$ is cotorsion
for every index\/~$\alpha$. \par
\textup{(b)} The $R$\+module $C$ is injective if and only if
the $S_\alpha$\+module\/ $\Hom_R(S_\alpha,C)$ is injective
for every index\/~$\alpha$.
\end{lem}

\begin{proof}
 Part~(a) is~\cite[Lemma~1.3.6(a)]{Pcosh} or~\cite[Example~3.8]{Pal},
while part~(b) is~\cite[Lemma~1.3.6(b)]{Pcosh}
or~\cite[Example~3.7]{Pal}.
 The ``only if'' implication in part~(a) holds by
Lemma~\ref{cotorsion-change-of-scalars-lemma}(b),
and the ``only if'' implication in part~(a) holds by
Lemma~\ref{injective-change-of-scalars-lemma}(b).
 The ``if'' implications in both parts~(a) and~(b) are provable
using Lemma~\ref{contraadjusted-cech-resolution}.

 Let us sketch the proof of the ``if'' implication in part~(a).
 Using Lemma~\ref{cotorsion-change-of-scalars-lemma}(a\+-b),
one shows that all the terms of the exact
sequence~\eqref{cta-module-cech-resolution}, except perhaps
the rightmost one, are cotorsion $R$\+modules.
 Then one proceeds by induction, moving from the leftmost to
the rightmost end of the sequence~\eqref{cta-module-cech-resolution}
and proving that its $R$\+modules of cycles are cotorsion.
 The fact that the class of cotorsion $R$\+modules is closed under
cokernels of monomorphisms in $R\Modl$ (see
Section~\ref{cotorsion-modules-subsecn}) plays a key role in
this argument.
\end{proof}

\begin{lem} \label{short-exact-sequnces-of-cta-are-colocal}
 Let $C\rarrow D\rarrow E$ be a composable pair of homomorphisms of
contraadjusted $R$\+modules.
 Then\/ $0\rarrow C\rarrow D\rarrow E\rarrow0$ is a short exact
sequence of $R$\+modules if and only if\/ $0\rarrow\Hom_R(S_\alpha,C)
\rarrow\Hom_R(S_\alpha,D)\rarrow\Hom_R(S_\alpha,E)\rarrow0$ is a short
exact sequence of $S_\alpha$\+modules for every index\/~$\alpha$.
\end{lem}

\begin{proof}
 This is~\cite[Lemma~1.4.1(a)]{Pcosh}.
 In particular, the easy ``only if'' assertion follows from
Example~\ref{open-affine-very-flat-example} above.
\end{proof}

\begin{lem} \label{colocality-of-epimorphisms}
\textup{(a)} Let $D\rarrow E$ be a homomorphism of contraadjusted
$R$\+modules.
 Then $D\rarrow E$ is an admissible epimorphism in $R\Modl^\cta$ if
and only if\/ $\Hom_R(S_\alpha,D)\rarrow\Hom_R(S_\alpha,E)$ is
an admissible epimorphism in $S_\alpha\Modl^\cta$ for every
index\/~$\alpha$. {\hbadness=1350\par}
\textup{(b)} Let $D\rarrow E$ be a homomorphism of cotorsion
$R$\+modules.
 Then $D\rarrow E$ is an admissible epimorphism in $R\Modl^\cot$ if
and only if\/ $\Hom_R(S_\alpha,D)\rarrow\Hom_R(S_\alpha,E)$ is
an admissible epimorphism in $S_\alpha\Modl^\cot$ for every
index\/~$\alpha$. \par
\textup{(c)} Let $D\rarrow E$ be a homomorphism of injective
$R$\+modules.
 Then $D\rarrow E$ is a split epimorphism if and only if\/
$\Hom_R(S_\alpha,D)\rarrow\Hom_R(S_\alpha,E)$ is a split
epimorphism of (injective) $S_\alpha$\+modules for every
index\/~$\alpha$.
\end{lem}

\begin{proof}
 This is~\cite[Lemmas~1.4.1(b), 1.4.2(b), and~1.4.3(b)]{Pcosh}.
\end{proof}


\Section{Locally Contraherent Cosheaves}

\subsection{Copresheaves and cosheaves}
 Let $(X,\cO_X)$ be a ringed space.
 We skip the well-known definitions of presheaves and sheaves of
$\cO_X$\+modules on~$X$.
 The notation $(X,\cO_X)\Sh$ stands for the Grothendieck abelian
category of sheaves of $\cO_X$\+modules on~$X$.

 Let $\bB$ be a set of open subsets forming a base of the topology
of~$X$.
 We will view $\bB$ as a category: the objects are the open subsets
of $X$ belonging to $\bB$, and the morphisms are the identity
inclusions.

 By a \emph{presheaf of abelian groups on\/~$\bB$} one means
a contravariant functor $\bB^\sop\rarrow\Ab$ (where $\Ab=\boZ\Modl$
is the category of abelian groups).
 Given two open subsets $V\subset U\subset X$, \ $U$, $V\in\bB$,
and a presheaf $\cM$ on $\bB$, the homomorphism of abelian groups
$\cM(U)\rarrow\cM(V)$ assigned to the identity inclusion $V\rarrow U$
by the presheaf $\cM$ is called the \emph{restriction map}.

 A \emph{presheaf of\/ $\cO_X$\+modules $\cM$ on\/~$\bB$} is a presheaf
of abelian groups such that, for every open subset $U\subset X$,
\ $U\in\bB$, the abelian group $\cM(U)$ is endowed with
an $\cO_X(U)$\+module structure, and for every pair of open subsets
$V\subset U\subset X$, \ $U$, $V\in\bB$, the restriction map
$\cM(U)\rarrow\cM(V)$ is a homomorphism of $\cO_X(U)$\+modules.
 Here the $\cO_X(V)$\+module $\cM(V)$ is endowed with
a $\cO_X(U)$\+module structure using the restriction homomorphism
of rings $\cO_X(U)\rarrow\cO_X(V)$.

 A presheaf of abelian groups $\cM$ on $\bB$ is called a \emph{sheaf
of abelian groups} (\emph{on\/~$\bB$}) if the following \emph{sheaf
axiom} is satisfied.
 Let $U\subset X$, \ $U\in\bB$, be an open subset, and let
$U=\bigcup_\alpha V_\alpha$ be an open covering of $U$ such that
$V_\alpha\in\bB$ for all the indices~$\alpha$.
 Furthermore, for every pair of indices $\alpha$ and~$\beta$, let
$V_\alpha\cap V_\beta=\bigcup_\gamma W_{\alpha,\beta,\gamma}$ be
an open covering of the intersection $V_\alpha\cap V_\beta$
such that $W_{\alpha\beta\gamma}\in\bB$ for all~$\gamma$.
 Then the sequence of abelian groups
\begin{equation} \label{sheaf-axiom-for-topology-base}
 0\lrarrow \cM(U)\lrarrow\prod\nolimits_\alpha\cM(V_\alpha)
 \lrarrow\prod\nolimits_{\alpha,\beta,\gamma}\cM(W_{\alpha\beta\gamma})
\end{equation}
should be left exact.

 A presheaf of $\cO_X$\+modules $\cM$ on $\bB$ is called a \emph{sheaf
of\/ $\cO_X$\+modules} (\emph{on\/~$\bB$}) if its underlying presheaf
of abelian groups $\cM$ is a sheaf of abelian groups on~$\bB$.
 We denote the category of sheaves of $\cO_X$\+modules on $\bB$ by
$(\bB,\cO_X)\Sh$.

\begin{prop} \label{sheaves-determined-on-topology-base}
 The functor of restriction of sheaves of\/ $\cO_X$\+modules on $X$
to a topology base\/ $\bB$ is an equivalence of categories
$$
 (X,\cO_X)\Sh\simeq(\bB,\cO_X)\Sh.
$$
 In other words, any sheaf of\/ $\cO_X$\+modules on\/ $\bB$ can be
extended to a sheaf of\/ $\cO_X$\+modules on the whole of~$X$ (i.~e.,
on all the open subsets of~$X$) in a unique way, and any morphism
of sheaves of\/ $\cO_X$\+modules on\/ $\bB$ can be uniquely extended
to a morphism of sheaves of\/ $\cO_X$\+modules on~$X$.
\end{prop}

\begin{proof}
 This is~\cite[Section~0.3.2]{EGAI}, \cite[Lemma Tag~009U]{SP},
\cite[Proposition~2.1.3]{Pcosh}, \cite[Theorem~2.1(a)]{Pdomc},
or~\cite[Section~3.2]{Pform}.
\end{proof}

 Now we turn to copresheaves and cosheaves.
 A \emph{copresheaf of abelian groups} on\/ $\bB$ is a covariant
functor $\bB\rarrow\Ab$.
 We denote the abelian group assigned by a copresheaf $\fP$ to
an open subset $U\subset X$, \ $U\in\bB$, by~$\fP[U]$.
 For any two open subsets $V\subset U\subset X$, \ $U$, $V\in\bB$,
and a copresheaf $\fP$ on $\bB$, the homomorphism of abelian groups
$\fP[V]\rarrow\fP[U]$ assigned to the identity inclusion $V\rarrow U$
by the copresheaf $\fP$ is called the \emph{corestriction map}.

 A \emph{copresheaf of\/ $\cO_X$\+modules\/ $\fP$ on\/~$\bB$} is
a copresheaf of abelian groups such that, for every open subset
$U\subset X$, \ $U\in\bB$, the abelian group $\fP[U]$ is endowed with
an $\cO_X(U)$\+module structure, and for every pair of open subsets
$V\subset U\subset X$, \ $U$, $V\in\bB$, the corestriction map
$\fP[V]\rarrow\fP[U]$ is a homomorphism of $\cO_X(U)$\+modules.
 Here, as above, the $\cO_X(V)$\+module $\fP[V]$ is endowed with
a $\cO_X(U)$\+module structure using the restriction homomorphism
of rings $\cO_X(U)\rarrow\cO_X(V)$.

 A copresheaf of abelian groups $\fP$ on $\bB$ is called a \emph{cosheaf
of abelian groups} (\emph{on\/~$\bB$}) if the following \emph{cosheaf
axiom} is satisfied.
 As above, we consider an open subset $U\subset X$, \ $U\in\bB$
and an open covering $U=\bigcup_\alpha V_\alpha$ such that
$V_\alpha\in\bB$ for all the indices~$\alpha$.
 Furthermore, for every pair of indices $\alpha$ and~$\beta$,
we choose an open covering $V_\alpha\cap V_\beta=
\bigcup_\gamma W_{\alpha,\beta,\gamma}$ of the intersection
$V_\alpha\cap V_\beta$ such that $W_{\alpha\beta\gamma}\in\bB$
for all~$\gamma$.
 For any such choices of $U$, $V_\alpha$, and $W_{\alpha\beta\gamma}$,
the sequence of abelian groups
\begin{equation} \label{cosheaf-axiom-for-topology-base}
 \bigoplus\nolimits_{\alpha,\beta,\gamma}\fP[W_{\alpha\beta\gamma}]
 \lrarrow\bigoplus\nolimits_\alpha\fP[V_\alpha]\lrarrow\fP[U]\lrarrow0
\end{equation}
should be right exact.

 A copresheaf of $\cO_X$\+modules $\fP$ on $\bB$ is called
a \emph{cosheaf of\/ $\cO_X$\+modules} (\emph{on\/~$\bB$}) if its
underlying copresheaf of abelian groups $\fP$ is a cosheaf of abelian
groups on~$\bB$.
 We denote the category of cosheaves of $\cO_X$\+modules on $\bB$ by
$(\bB,\cO_X)\Cosh$.

 A \emph{co}(\emph{pre})\emph{sheaf} (of abelian groups or
$\cO_X$\+modules) \emph{on~$X$} is a co(pre)sheaf on the topology base
consisting of all the open subsets of~$X$.
 The category of cosheaves of $\cO_X$\+modules on $X$ will be denoted
by $(X,\cO_X)\Cosh$.

\begin{prop} \label{cosheaves-determined-on-topology-base}
 The functor of restriction of cosheaves of\/ $\cO_X$\+modules on $X$
to a topology base\/ $\bB$ is an equivalence of categories
$$
 (X,\cO_X)\Cosh\simeq(\bB,\cO_X)\Cosh.
$$
 In other words, any cosheaf of\/ $\cO_X$\+modules on\/ $\bB$ can be
extended to a cosheaf of\/ $\cO_X$\+modules on the whole of~$X$ (i.~e.,
on all the open subsets of~$X$) in a unique way, and any morphism of
cosheaves of\/ $\cO_X$\+modules on\/ $\bB$ can be uniquely extended
to a morphism of cosheaves of\/ $\cO_X$\+modules on~$X$.
\end{prop}

\begin{proof}
 This is~\cite[Theorem~2.1.2]{Pcosh}, \cite[Theorem~2.1(b)]{Pdomc},
or~\cite[Section~3.2]{Pform}.
\end{proof}

\begin{rem} \label{one-covering-of-the-intersections-is-enough}
 The sheaf axiom~\eqref{sheaf-axiom-for-topology-base} implies, in
particular, that the map $\cM(U)\rarrow\prod_\alpha\cM(V_\alpha)$
is injective for any open covering $U=\bigcup_\alpha V_\alpha$ with
$U$, $V_\alpha\in\bB$.
 Using this fact, one can show that, once the open covering
$U=\bigcup_\alpha V_\alpha$ is fixed, it is not necessary to check
 the sheaf axiom for \emph{all} coverings $V_\alpha\cap V_\beta=
\bigcup_\gamma W_{\alpha\beta\gamma}$ of the intersections
$V_\alpha\cap V_\beta$ by open subsets $W_{\alpha\beta\gamma}\in\bB$.
 Chosing \emph{one} such covering for every pair of indices $\alpha$,
$\beta$ and checking~\eqref{sheaf-axiom-for-topology-base} for such
chosen coverings of the intersections $V_\alpha\cap V_\beta$ is
enough to establish the validity of the sheaf axiom for~$\cM$.

 Similarly, the cosheaf axiom~\eqref{cosheaf-axiom-for-topology-base}
implies, in particular, that the map $\bigoplus_\alpha\fP[V_\alpha]
\rarrow\fP[U]$ is surjective for any open covering
$U=\bigcup_\alpha V_\alpha$ with $U$, $V_\alpha\in\bB$.
 Using this fact, one can show that it is not necessary to check
the cosheaf axiom for \emph{all} coverings $V_\alpha\cap V_\beta=
\bigcup_\gamma W_{\alpha\beta\gamma}$ of the intersections
$V_\alpha\cap V_\beta$ by open subsets $W_{\alpha\beta\gamma}\in\bB$.
 Chosing \emph{one} such covering for every pair of indices $\alpha$,
$\beta$ and checking~\eqref{cosheaf-axiom-for-topology-base} for such
chosen coverings of the intersections $V_\alpha\cap V_\beta$ is
sufficient.
\end{rem}

\begin{rems} \label{finite-coverings-suffice-remarks}
 Cosheaves are dual-analogous to sheaves (cf.~\cite[Preface]{Pcosh}).
 In what ways can one make this informal statement precise?

\smallskip
 (a)~Let $\fP$ be a cosheaf of $\cO_X$\+modules on $X$ and $A$
be an arbitrary abelian group.
 Comparing the axioms~\eqref{sheaf-axiom-for-topology-base}
and~\eqref{cosheaf-axiom-for-topology-base}, one can see that the rule
$\cM(U)=\Hom_\boZ(\fP[U],A)$ for all open subsets $U\subset X$
defines a cosheaf of $\cO_X$\+modules on~$X$.
 In particular, one can define the \emph{character sheaf} $\fP^+$ of
a cosheaf $\fP$ by the rule $\fP^+(U)=\Hom_\boZ(\fP[U],\boQ/\boZ)$
for all open subsets $U\subset X$.

\smallskip
 (b)~Assume that the topological space $X$ has a topology base $\bB$
with the following properties:
\begin{itemize}
\item all open subsets $U\subset X$, \ $U\in\bB$ are quasi-compact
(in the topology on $U$ induced from the topology of~$X$);
\item for any open subset $U\subset X$, \ $U\in\bB$, and any two open
subsets $V$, $W\subset U$, \ $V$, $W\in\bB$, the intersection $V\cap W$
is quasi-compact.
\end{itemize}
 Then, for (co)presheaves on $\bB$, it suffices to check both the sheaf
axiom~\eqref{sheaf-axiom-for-topology-base} and the cosheaf
axiom~\eqref{cosheaf-axiom-for-topology-base} for \emph{finite} open
coverings $U=\bigcup_\alpha V_\alpha$ and $V_\alpha\cap V_\beta=
\bigcup_\gamma W_{\alpha\beta\gamma}$
\emph{only}~\cite[Remark~2.1.4]{Pcosh}, \cite[Remarks~3.2.1]{Pform}.

\smallskip
 (c)~In particular, let $X$ be a scheme.
 Then the topology base $\bB$ consisting of all the quasi-compact,
quasi-separated open subsets $U\subset X$ satisfies both the conditions
of~(b).
 Alternatively, the topology base consisting of all the affine open
subschemes $U\subset X$ also satisfies both the conditions of~(b).

\smallskip
 (d)~Under the assumptions of~(b), let $\cM$ be a sheaf of
$\cO_X$\+modules on $X$ and $J$ be an injective abelian group.
 Comparing the axioms~\eqref{sheaf-axiom-for-topology-base}
and~\eqref{cosheaf-axiom-for-topology-base} for \emph{finite} open
coverings, one can see that the rule $\fP[U]=\Hom_\boZ(\cM(U),J)$ for
all open subsets $U\subset X$, \ $U\in\bB$ defines a cosheaf of
$\cO_X$\+modules on~$\bB$.
 Using Proposition~\ref{cosheaves-determined-on-topology-base}, one
can uniquely extend the resulting cosheaf of $\cO_X$\+modules on $\bB$
to a cosheaf of $\cO_X$\+modules on~$X$.
 In particular, one can define the \emph{character cosheaf} $\cM^+$ of
a sheaf $\cM$ by the rule $\cM^+[U]=\Hom_\boZ(\cM(U),\boQ/\boZ)$
for all open subsets $U\subset X$, \ $U\in\bB$.

\smallskip
 (e)~It is also clear from the sheaf
axiom~\eqref{sheaf-axiom-for-topology-base} that infinite products
exist in the category $(X,\cO_X)\Sh$, and the functors $\cM\longmapsto
\cM(U)\:(X,\cO_X)\Sh\rarrow\cO_X(U)\Modl$ of sections over all open
subsets $U\subset X$ preserve infinite products.
 Similarly, it is clear from the cosheaf 
axiom~\eqref{cosheaf-axiom-for-topology-base} that infinite coproducts
exist in the category $(X,\cO_X)\Cosh$, and the functors $\fP\longmapsto
\fP[U]\:(X,\cO_X)\Cosh\rarrow\cO_X(U)\Modl$ of cosections over all open
subsets $U\subset X$ preserve infinite coproducts.

\smallskip
 (f)~Now, under the assumptions of~(b), it is clear from the sheaf
axiom~\eqref{sheaf-axiom-for-topology-base} \emph{for finite open
coverings} that infinite coproducts exist in the category
$(\bB,\cO_X)\Sh$, and the functors of sections over all open
subsets subsets $U\subset X$, \ $U\in\bB$ preserve infinite coproducts.
 In view of
Proposition~\ref{sheaves-determined-on-topology-base}, it follows
that infinite coproducts exist in the category $(X,\cO_X)\Sh$,
and the functors of sections over open subsets $U\subset X$, \ $U\in\bB$
preserve infinite coproducts.
 (These are well-known facts for sheaves, of course.)

 Similarly, under the assumptions of~(b), it is clear from the cosheaf
axiom~\eqref{cosheaf-axiom-for-topology-base} \emph{for finite open
coverings} that infinite products exist in the category
$(\bB,\cO_X)\Cosh$, and the functors of cosections over all open
subsets subsets $U\subset X$, \ $U\in\bB$ preserve infinite products.
 In view of
Proposition~\ref{cosheaves-determined-on-topology-base}, it follows
that infinite products exist in the category $(X,\cO_X)\Cosh$,
and the functors of cosections over open subsets $U\subset X$, \
$U\in\bB$ preserve infinite products.
\end{rems}

\subsection{The contraadjustedness and contraherence conditions}
\label{contraadj-contraher-subsecn}
 The definitions of a \emph{contraherent cosheaf} and a \emph{locally
contraherent cosheaf} on a scheme are dual analogues of the definition
of quasi-coherent sheaves suggested by Enochs and Estrada
in~\cite{EE}.
 We start with spelling out the relevant formulation of quasi-coherence
before passing to contraherence.

 Let $X$ be a scheme with a chosen open covering~$\bW$.
 We will say that an open subscheme $U\subset X$ is \emph{subordinate
to\/~$\bW$} if there is an open subset $W\in\bW$ such that $U\subset W$.
 Denote by $\bB$ the topology base in $X$ consisting of all the affine
open subschemes $U\subset X$ subordinate to~$\bW$.

 Let $\cM$ be a presheaf of $\cO_X$\+modules on~$\bB$.
 The presheaf $\cM$ is said to be \emph{quasi-coherent} if
the following \emph{quasi-coherence axiom} is satisfied:
\begin{enumerate}
\renewcommand{\theenumi}{\roman{enumi}}
\item for any pair of affine open subschemes $U\subset V$ subordinate
to $\bW$ in $X$, the restriction map of $\cO_X(U)$\+modules
$\cM(U)\rarrow\cM(V)$ induces an isomorphism of $\cO_X(V)$\+modules
$$
 \cO_X(V)\ot_{\cO_X(U)}\cM(U)\simeq\cM(V).
$$
\end{enumerate}

 Any quasi-coherent presheaf $\cM$ on $\bB$ is a sheaf on~$\bB$ (as one
can see from the exactness of the leftmost fragment of the \v Cech sequence~\eqref{module-cech-coresolution}
together with Remarks~\ref{one-covering-of-the-intersections-is-enough}
and~\ref{finite-coverings-suffice-remarks}(b\+-c)).
 So, by Proposition~\ref{sheaves-determined-on-topology-base}, one
can uniquely extend $\cM$ to a sheaf of $\cO_X$\+modules on the whole
of~$X$ (i.~e., on all the open subsets of~$X$).

 Furthermore, the quasi-coherence property of sheaves of
$\cO_X$\+modules is \emph{local} in the following sense.
 Denote by $\bW_X=\{X\}$ the trivial open covering of $X$ and by
$\bB_X$ the topology base consisting of all the affine open
subschemes of~$X$.

 Let $\cM$ be a sheaf of $\cO_X$\+modules on~$X$.
 Then the restriction of $\cM$ to $\bB_X$ is quasi-coherent if and
and only if the restriction of $\cM$ to $\bB$ is quasi-coherent.
 In other words, the quasi-coherence axiom~(i) holds for all pairs
of affine open subschemes $V\subset U$ subordinate to $\bW$ in $X$
if and only if it holds for all pairs of affine open subschemes
$V\subset U$ in~$X$.
 (This is essentially a restatement of the classical theory of
quasi-coherent sheaves on affine schemes~\cite[Lemma Tag~01IA]{SP};
cf.~\cite[Proposition~3.4.2]{Pform}.)
 If this is the case, the sheaf of $\cO_X$\+modules $\cM$ is called
\emph{quasi-coherent}.

 We denote the category of quasi-coherent sheaves on $X$ by $X\Qcoh$.
 It is a Grothendieck abelian category~\cite[Proposition Tag~077P]{SP}.

 Now let $\fP$ be a copresheaf of $\cO_X$\+modules on~$\bB$.
 The copresheaf $\fP$ is said to be \emph{contraherent} if
the following two axioms are satisfied:
\begin{enumerate}
\renewcommand{\theenumi}{\roman{enumi}}
\setcounter{enumi}{1}
\item for any pair of affine open subschemes $U\subset V$ subordinate
to $\bW$ in $X$, the corestriction map of $\cO_X(U)$\+modules
$\fP[V]\rarrow\fP[U]$ induces an isomorphism of $\cO_X(V)$\+modules
$$
 \fP[V]\simeq\Hom_{\cO_X(U)}(\cO_X(V),\fP[U]);
$$
\item for any pair of affine open subschemes $U\subset V$ subordinate
to $\bW$ in $X$, one has
$$
 \Ext^1_{\cO_X(U)}(\cO_X(V),\fP[U])=0.
$$
\end{enumerate}
 The projective dimension of the $\cO_X(U)$\+module $\cO_X(V)$ never
exceeds~$1$ \,\cite[Remark~5.1]{Pphil}; that is why no $\Ext^n$
vanishing condition for $n\ge2$ is needed in~(iii).

 Fix an affine open subscheme $U\subset X$ subordinate to~$\bW$.
 Then condition~(iii) holds for all affine open subschemes $V\subset U$
if and only if the $\cO_X(U)$\+module $\fP[U]$ is contraadjusted.
 This is clear from Example~\ref{open-affine-very-flat-example}.
 For this reason, condition~(iii) is called
the \emph{contraadjustedness axiom}.
 Condition~(ii) is called the \emph{contraherence axiom}.

 Any contraherent copresheaf $\fP$ on $\bB$ is a cosheaf on~$\bB$ (as
one can see from the exactness of the rightmost fragment of the \v Cech sequence~\eqref{cta-module-cech-resolution}
together with Remarks~\ref{one-covering-of-the-intersections-is-enough}
and~\ref{finite-coverings-suffice-remarks}(b\+-c)).
 Hence, by Proposition~\ref{cosheaves-determined-on-topology-base}, one
can uniquely extend $\fP$ to a cosheaf of $\cO_X$\+modules on the whole
of~$X$ (i.~e., on all the open subsets of~$X$).

 The contraadjustedness condition for cosheaves of $\cO_X$\+modules
on $X$ is local, but the contaherence condition is \emph{not local}.
 A cosheaf of $\cO_X$\+modules $\fP$ on $X$ is said to be
\emph{$\bW$\+locally contraherent} if the restriction of $\fP$ to
the topology base $\bB$ is contraherent, i.~e., if
the contraadjustedness and contraherence axioms~(ii\+-iii) hold
for all pairs of affine open subschemes $V\subset U$ subordinate
to $\bW$ in~$X$.
 A cosheaf of $\cO_X$\+modules $\fP$ is called \emph{locally
contraherent} if it is $\bW$\+locally contraherent for \emph{some}
open covering $\bW$ of~$X$.

 A cosheaf of $\cO_X$\+modules $\fP$ on $X$ is called
\emph{contraherent} if it is $\bW_X$\+locally contraherent for
the trivial open covering $\bW_X=\{X\}$, i.~e., if the restriction
of $\fP$ to the topology base $\bB_X$ is contraherent.
 In other words, a cosheaf $\fP$ on $X$ is contraherent if and only if
the contraadjustedness and contraherence axioms~(ii\+-iii) hold
for all pairs of affine open subschemes $V\subset U$ in~$X$.

 A counterexample of a locally contraherent, but not contraherent
cosheaf of $\cO_X$\+modules on a scheme $X$ can be found
in~\cite[Example~3.2.1]{Pcosh}.
 This counterexample is quite simple in nature; taking $X$ to be
the spectrum of the ring of integers $\boZ$ or the ring of polynomials
in one variable $k[x]$ over a field~$k$ is enough.
 The locality of the contraadjustedness condition is explained
in~\cite[paragraph after Example~3.2.1]{Pcosh}.

 We denote the additive category of $\bW$\+locally contraherent
cosheaves on $X$ by $X\Lcth_\bW$.
 The additive category of (globally) contraherent cosheaves is
denoted by $X\Ctrh=X\Lcth_{\{X\}}$.
 The additive category of all locally contraherent cosheaves on $X$
is denoted by $X\Lcth=\bigcup_\bW X\Lcth_\bW$.
 So we have inclusions of full subcategories
$$
 X\Ctrh\subset X\Lcth_\bW\subset X\Lcth.
$$

\subsection{Locally contraherent cosheaves}
 The categories of (locally) contraherent cosheaves are almost never
abelian; but they always carry natural exact category structures.
 Let us define these exact structures.

 Let $X$ be a scheme with an open covering~$\bW$.
 A composable pair of morphisms of $\bW$\+locally contraherent cosheaves
$\fP\rarrow\fQ\rarrow\fR$ is said to be an (\emph{admissible})
\emph{short exact sequence} $0\rarrow\fP\rarrow\fQ\rarrow\fR\rarrow0$
in $X\Lcth_\bW$ if, for every affine open subscheme $U\subset X$
subordinate to $\bW$, the short sequence of $\cO_X(U)$\+modules
$0\rarrow\fP[U]\rarrow\fQ[U]\rarrow\fR[U]\rarrow0$ is exact.

 In view of Lemma~\ref{short-exact-sequnces-of-cta-are-colocal}, it
suffices to check this condition for affine open subschemes $U$
belonging to any chosen affine open covering of the scheme $X$
subordinate to the open covering~$\bW$.
 In this sense, the property of a short sequence of $\bW$\+locally
contraherent cosheaves on $X$ to be exact is \emph{local}.
 Using Lemma~\ref{Hom-into-pushout-lemma}(b) together with
Example~\ref{open-affine-very-flat-example}, one can easily show
that the class of short exact sequences defined above is indeed
an exact category structure on $X\Lcth_\bW$.

 Let $\bW'$ be another open covering of~$X$.
 We say that $\bW'$ is \emph{subordinate to\/~$\bW$} if every open
subscheme belonging to $\bW'$ is subordinate to $\bW$ (in the sense
of the definition in Section~\ref{contraadj-contraher-subsecn}).
 Given an open covering $\bW'$ subordinate to $\bW$, it follows
from the previous paragraph that a short sequence in $X\Lcth_\bW$
is exact in $X\Lcth_\bW$ if and only if it is exact as a short
sequence in $X\Lcth_{\bW'}$.

 It follows from Lemma~\ref{colocality-of-epimorphisms}(a) that
the full subcategory $X\Lcth_\bW$ is closed under kernels of
admissible epimorphisms in $X\Lcth_{\bW'}$.
 So a morphism in $X\Lcth_\bW$ is an admissible epimorphism in
$X\Lcth_\bW$ if and only if it is an admissible epimorphism in
$X\Lcth_{\bW'}$.
 Furthermore, the homological criterion of contraherence of
a locally contraherent cosheaf on an affine
scheme~\cite[Lemma~3.2.2]{Pcosh} (see also~\cite[Corollary~4.6.1]{Pcosh}
and/or~\cite[Proposition~3.6.1]{Pform}) implies that the full
subcategory $X\Lcth_\bW$ is closed under extensions in $X\Lcth_{\bW'}$.

 Let us emphasize, however, that the full subcategory $X\Lcth_\bW$
is usually \emph{not} closed under cokernels of admissible
monomorphisms in $X\Lcth_{\bW'}$.
 See~\cite[Example~3.2.1]{Pcosh}.

 A short sequence in the category $X\Lcth$ is called \emph{exact}
if it is a short exact sequence in $X\Lcth_\bW$ for some open
covering $\bW$ of~$X$.
 So the full subcategory $X\Lcth_\bW$ is closed under extensions
and kernels of admissible epimorphisms in the exact category
$X\Lcth$, and a short sequence in $X\Lcth_\bW$ is exact if and only
if it is exact in $X\Lcth$.

 Specializing the discussion above in this section to the case of
the trivial covering $\bW_X=\{X\}$, we obtain an exact category
structure on the category $X\Ctrh=X\Lcth_{\bW_X}$ of contraherent
cosheaves on~$X$.

\begin{lem} \label{ctrh-cosheaves-on-affine-scheme}
 Let $U$ be an affine scheme.
 Then the functor\/ $\fP\longmapsto\fP[U]$ establishes an exact
category equivalence $U\Ctrh\simeq\cO(U)\Modl^\cta$ between
the exact category of contraherent cosheaves on $U$ and the exact
category of contraadjusted modules over the ring\/~$\cO(U)$.
\end{lem}

\begin{proof}
 Follows immediately from the definitions and the discussion above.
\end{proof}

\subsection{The local cotorsion and local injectivity conditions}
\label{lct-lin-subsecn}
 Let $X$ be a scheme with an open covering~$\bW$.
 A $\bW$\+locally contraherent cosheaf $\fP$ on $X$ is said to be
\emph{locally cotorsion} if the $\cO_X(U)$\+module $\fP[U]$ is
cotorsion for every affine open subscheme $U\subset X$ subordinate
to~$\bW$.
 In view of
Lemma~\ref{cotorsion-injective-colocal-if-contraadjusted}(a), it
suffices to check this condition for affine open subschemes $U$
belonging to any chosen affine open covering of~$X$.

 In particular, let $\bW'$ be an affine open covering of $X$
subordinate to~$\bW$.
 Then it follows from the previous paragraph that a $\bW$\+locally
contraherent cosheaf $\fP$ on $X$ is locally cotorsion if and only
if $\fP$ is a locally cotorsion $\bW'$\+locally contraherent
cosheaf on~$X$.

 We denote the full subcategory of locally cotorsion $\bW$\+locally
contraherent cosheaves on $X$ by $X\Lcth^\lct_\bW\subset X\Lcth_\bW$.
 The notation $X\Lcth^\lct=\bigcup_\bW X\Lcth^\lct_\bW\allowbreak
\subset X\Lcth$ stands for the full subcategory of locally cotorsion
locally contraherent cosheaves in $X\Lcth$.

 In particular, a contraherent cosheaf $\fP$ on $X$ is said to be
\emph{locally cotorsion} if it belongs to the full subcategory
$X\Ctrh^\lct=X\Lcth^\lct_{\{X\}}\subset X\Ctrh$.
 So we have inclusions of full subcategories
$$
  X\Ctrh^\lct\subset X\Lcth^\lct_\bW\subset X\Lcth^\lct.
$$

 Clearly, the full subcategory $X\Lcth^\lct_\bW$ is closed under
extensions and cokernels of admissible monomorphisms in $X\Lcth_\bW$.
 So the additive category $X\Lcth^\lct_\bW$ carries the exact
category structure inherited from $X\Lcth_\bW$.
 In particular, in the case of the trivial open covering
$\bW_X=\{X\}$, we obtain an exact structure on the category of
locally cotorsion contraherent cosheaves $X\Ctrh^\lct$.

 Similarly, the full subcategory $X\Lcth^\lct$ is closed under
extensions and cokernels of admissible monomorphisms in $X\Lcth$,
so we have the inherited exact category structure on $X\Lcth^\lct$.
 It follows that the full exact subcategories $X\Ctrh^\lct$, \
$X\Lcth^\lct_\bW$, and $X\Lcth^\lct$ are closed under extensions
and kernels of admissible epimorphisms in each other (see
also Lemma~\ref{colocality-of-epimorphisms}(b)).

 A $\bW$\+locally contraherent cosheaf $\fJ$ on $X$ is said to be
\emph{locally injective} if the $\cO_X(U)$\+module $\fJ[U]$ is
injective for every affine open subscheme $U\subset X$ subordinate
to~$\bW$.
 In view of
Lemma~\ref{cotorsion-injective-colocal-if-contraadjusted}(b), it
suffices to check this condition for affine open subschemes $U$
belonging to any chosen affine open covering of~$X$.
 In particular, it follows that a $\bW$\+locally contraherent cosheaf
$\fJ$ on $X$ is locally injective if and only if $\fJ$ is a locally
injective $\bW'$\+locally contraherent cosheaf on~$X$.

 We denote the full subcategory of locally injective $\bW$\+locally
contraherent cosheaves on $X$ by $X\Lcth^\lin_\bW\subset X\Lcth_\bW$.
 The notation $X\Lcth^\lin=\bigcup_\bW X\Lcth^\lin_\bW\allowbreak
\subset X\Lcth$ stands for the full subcategory of locally injective
locally contraherent cosheaves in $X\Lcth$.  {\hbadness=1150\par}

 In particular, a contraherent cosheaf $\fJ$ on $X$ is said to be
\emph{locally injective} if it belongs to the full subcategory
$X\Ctrh^\lin=X\Lcth^\lin_{\{X\}}\subset X\Ctrh$.
 So we have inclusions of full subcategories
$$
  X\Ctrh^\lin\subset X\Lcth^\lin_\bW\subset X\Lcth^\lin,
$$
as well as
$$
 X\Ctrh^\lin\subset X\Ctrh^\lct\subset X\Ctrh, \qquad
 X\Lcth^\lin_\bW\subset X\Lcth^\lct_\bW\subset X\Lcth_\bW,
$$
and
$$
 X\Lcth^\lin\subset X\Lcth^\lct\subset X\Lcth.
$$

 Clearly, the full subcategory $X\Lcth^\lin_\bW$ is closed under
extensions and cokernels of admissible monomorphisms in $X\Lcth_\bW$
and in $X\Lcth_\bW^\lct$.
 So the additive category $X\Lcth^\lin_\bW$ carries the exact
category structure inherited from $X\Lcth_\bW$ or $X\Lcth^\lct_\bW$.
 In particular, in the case of the trivial open covering
$\bW_X=\{X\}$, we obtain an exact structure on the category of
locally injective contraherent cosheaves $X\Ctrh^\lin$.

 Similarly, the full subcategory $X\Lcth^\lin$ is closed under
extensions and cokernels of admissible monomorphisms in $X\Lcth$
and in $X\Lcth^\lct$, so we have the inherited exact category
structure on $X\Lcth^\lin$.
 It follows that the full exact subcategories $X\Ctrh^\lin$, \
$X\Lcth^\lin_\bW$, and $X\Lcth^\lin$ are closed under extensions
and kernels of admissible epimorphisms in each other (see
also Lemma~\ref{colocality-of-epimorphisms}(c)).

\begin{lem} \label{ctrh-lct-lin-cosheaves-on-affine-scheme}
 Let $U$ be an affine scheme.
 Then the equivalence of exact categories $U\Ctrh\simeq\cO(U)\Modl^\cta$
from Lemma~\ref{ctrh-cosheaves-on-affine-scheme} restricts to \par
\textup{(a)} an equivalence $U\Ctrh^\lct\simeq\cO(U)\Modl^\cot$ between 
the exact category of locally cotorsion contraherent cosheaves on $U$
and the exact category of cotorsion\/ $\cO(U)$\+modules; \par
\textup{(b)} an equivalence $U\Ctrh^\lin\simeq\cO(U)\Modl^\inj$ between
the split exact category of locally injective contraherent cosheaves
on $U$ and the additive (split exact) category of injective\/
$\cO(U)$\+modules. \qed
\end{lem}

\begin{rem} \label{character-cosheaf-ctrh-lct-remark}
 Let $X$ be a scheme and $\cM$ be a quasi-coherent sheaf on~$X$.
 Then the construction of
Remarks~\ref{finite-coverings-suffice-remarks}(b\+-d) provides
a cosheaf of $\cO_X$\+modules $\cM^+$ on $X$ defined by the rule
$\cM^+[U]=\Hom_\boZ(\cM(U),\boQ/\boZ)$ for all quasi-compact,
quasi-separated open subschemes $U\subset X$.
 Comparing the quasi-coherence axiom~(i) with the contraherence
axiom~(ii) from Section~\ref{contraadj-contraher-subsecn}, and taking
into account Lemma~\ref{flat-cotorsion-injective-tensor-Hom-lemma}(c),
one can see that $\cM^+$ is a locally cotorsion contraherent cosheaf
on $X$, that is, $\cM^+\in X\Ctrh^\lct$.
 For a flat quasi-coherent sheaf $\cF$ on $X$, the contraherent cosheaf
$\cF^+$ on $X$ is locally injective by
Lemma~\ref{flat-cotorsion-injective-tensor-Hom-lemma}(d).
 A further discussion of this construction can be found 
in~\cite[beginning of Section~2.4]{Pcosh}.
\end{rem}


\Section{Direct Images of Locally Contraherent Cosheaves}

\subsection{Direct image of cosheaves of modules}
 Let $f\:(Y,\cO_Y)\rarrow(X,\cO_X)$ be a morphism of ringed spaces.
 Given a presheaf of $\cO_Y$\+modules $\cN$ on $Y$, the presheaf
of $\cO_X$\+modules $f_*\cN$ on $X$ is defined by the rule
$$
 (f_*\cN)(U)=\cN(f^{-1}(U)),
$$
where $f^{-1}(U)\subset Y$ is the preimage of an open subset
$U\subset X$ under the continuous map of topological spaces
$f\:Y\rarrow X$.
 The $\cO_Y(f^{-1}(U))$\+module $\cN(f^{-1}(U))$ is endowed with
a $\cO_X(U)$\+module structure using the ring homomorphism
$\cO_X(U)\rarrow\cO_Y(f^{-1}(U))$ that is a part of the structure of
a morphism of ringed spaces $f\:(Y,\cO_Y)\rarrow(X,\cO_X)$.
{\hbadness=1250\par}

 For any open covering $U=\bigcup_\alpha V_\alpha$ of an open
subset $U\subset X$, the open subsets $f^{-1}(V_\alpha)\subset Y$
form an open covering $f^{-1}(U)=\bigcup_\alpha f^{-1}(V_\alpha)$
of the open subset $f^{-1}(U)\subset Y$.
 Using this observation, one can readily check that the presheaf
$f_*\cN$ on $X$ satisfies the sheaf
axiom~\eqref{sheaf-axiom-for-topology-base} whenever so does
the presheaf $\cN$ on~$Y$.
 So we have the functor of direct image of sheaves of $\cO$\+modules
\begin{equation} \label{direct-image-of-sheaves-of-O-modules}
 f_*\:(Y,\cO_Y)\Sh\lrarrow(X,\cO_X)\Sh.
\end{equation}

 Dual-analogously, given a copresheaf of $\cO_Y$\+modules $\fQ$ on $Y$,
the presheaf of $\cO_X$\+modules $f_!\fQ$ on $X$ is defined by
the rule
$$
 (f_!\fQ)[U]=\fQ[f^{-1}(U)]
$$
for all open subsets $U\subset X$.
 The $\cO_Y(f^{-1}(U))$\+module $\fQ[f^{-1}(U)]$ is endowed with
a $\cO_X(U)$\+module structure using the same ring homomorphism
$\cO_X(U)\rarrow\cO_Y(f^{-1}(U))$ as above.

 Using the same observation about the preimages of open coverings as
above, one can readily check that the copresheaf $f_!\fQ$ on $X$
satisfies the cosheaf axiom~\eqref{cosheaf-axiom-for-topology-base} 
whenever so does the copresheaf $\fQ$ on~$Y$.
 So we have the functor of direct image of cosheaves of $\cO$\+modules
\begin{equation} \label{direct-image-of-cosheaves-of-O-modules}
 f_!\:(Y,\cO_Y)\Cosh\lrarrow(X,\cO_X)\Cosh.
\end{equation}

\subsection{Direct image and contraherence}
\label{direct-image-and-contraherence-subsecn}
 Let $f\:Y\rarrow X$ be a morphism of schemes.
 Whenever the morphism~$f$ is quasi-compact and quasi-separated,
the direct image functor $f_*\:(Y,\cO_Y)\Sh\rarrow(X,\cO_X)\Sh$
\,\eqref{direct-image-of-sheaves-of-O-modules} takes quasi-coherent
sheaves on $Y$ to quasi-coherent sheaves on~$X$
\,\cite[Proposition~I.6.7.1]{EGAI}, \cite[Lemma Tag~01LC]{SP}
(cf.~\cite[Proposition~3.9.1(b)]{Pform}).

 The theory of direct images of contraherent and locally contraherent
cosheaves is a bit more complicated.
 Generally speaking, for an open immersion $f\:Y\rarrow X$ of smooth
algebraic varieties (i.~e., schemes of finite type) over a field,
even when the complement $X\setminus Y$ is smooth,
the direct image functor $f_!\:(Y,\cO_Y)\Cosh\rarrow(X,\cO_X)\Cosh$
\,\eqref{direct-image-of-cosheaves-of-O-modules} does not even take
contraherent cosheaves on $Y$ to locally contraherent cosheaves on~$X$
\,\cite[Remark~2.3.1 and Example~2.3.2]{Pcosh}.

 In order to make sure that the cosheaf $f_!\fQ$ is (locally)
contraherent, one needs to either assume the morphism~$f$ to be
affine (and more), or impose suitable adjustedness conditions
on a (locally) contraherent cosheaf~$\fQ$.
 In the rest of this section, we discuss affine morphisms~$f$.

 Firstly, let $f\:Y\rarrow X$ be an affine morphism of schemes.
 Then the direct image functor
$f_!\:(Y,\cO_Y)\Cosh\rarrow(X,\cO_X)\Cosh$
\,\eqref{direct-image-of-cosheaves-of-O-modules}
takes contraherent cosheaves on $Y$ to contraherent cosheaves on $X$,
and induces an exact functor between the respective exact categories
\begin{equation} \label{ctrh-direct-image}
 f_!\:Y\Ctrh\lrarrow X\Ctrh.
\end{equation}
 Furthermore, the same direct image functor
$f_!\:(Y,\cO_Y)\Cosh\rarrow(X,\cO_X)\Cosh$ takes locally cotorsion
contraherent cosheaves on $Y$ to locally cotorsion contraherent
cosheaves on $X$, and induces an exact functor between their exact
categories
\begin{equation} \label{ctrh-lct-direct-image}
 f_!\:Y\Ctrh^\lct\lrarrow X\Ctrh^\lct.
\end{equation}
 Finally, if the morphism~$f$ is affine \emph{and} flat,
then the same functor of direct image of cosheaves of
$\cO$\+modules~$f_!$ takes locally injective contraherent cosheaves
on $Y$ to locally injective contraherent cosheaves on $X$, and induces
an exact functor between the exact categories
\begin{equation} \label{ctrh-lin-direct-image}
 f_!\:Y\Ctrh^\lin\lrarrow X\Ctrh^\lin.
\end{equation}
 The proofs of these assertions use
Lemmas~\ref{contraadjusted-change-of-scalars-lemma}(a),
\ref{cotorsion-change-of-scalars-lemma}(a),
and~\ref{injective-change-of-scalars-lemma}(a).
 We refer to~\cite[Section~2.3]{Pcosh} for the details
(see also~\cite[Section~2.4]{Pdomc}).

 Secondly, let $\bW$ be an open covering of the scheme $X$ and $\bT$
be an open covering of the scheme~$Y$.
 One says that a morphism of schemes $f\:Y\rarrow X$ is
\emph{$(\bW,\bT)$\+affine} if, for every affine open subscheme
$U\subset X$ subordinate to $\bW$, the preimage $f^{-1}(U)\subset Y$
is an affine open subscheme subordinate to~$\bT$.
 Every $(\bW,\bT)$\+affine morphism of schemes is
affine~\cite[D\'efinition~II.1.2.1 and Corollaire~II.1.3.2]{EGAII},
\cite[Lemma Tag~01S8]{SP}.

 Let $f\:Y\rarrow X$ be a $(\bW,\bT)$\+affine morphism of schemes.
 Then the direct image functor
$f_!\:(Y,\cO_Y)\Cosh\rarrow(X,\cO_X)\Cosh$
\,\eqref{direct-image-of-cosheaves-of-O-modules}
takes $\bT$\+locally contraherent cosheaves on $Y$ to $\bW$\+locally
contraherent cosheaves on $X$, and induces an exact functor between
the respective exact categories
\begin{equation} \label{lcth-W-T-direct-image}
 f_!\:Y\Lcth_\bT\lrarrow X\Lcth_\bW.
\end{equation}
 Furthermore, the same direct image functor
$f_!\:(Y,\cO_Y)\Cosh\rarrow(X,\cO_X)\Cosh$ takes locally cotorsion
$\bT$\+locally contraherent cosheaves on $Y$ to locally cotorsion
$\bW$\+locally contraherent cosheaves on $X$, and induces an exact
functor between their exact categories
\begin{equation} \label{lcth-W-T-lct-direct-image}
 f_!\:Y\Lcth^\lct_\bT\lrarrow X\Lcth^\lct_\bW.
\end{equation}
 Finally, if the morphism~$f$ is $(\bW,\bT)$\+affine \emph{and} flat,
then the same functor of direct image of cosheaves of
$\cO$\+modules~$f_!$ takes locally injective $\bT$\+locally
contraherent cosheaves on $Y$ to locally injective $\bW$\+locally
contraherent cosheaves on $X$, and induces an exact functor
between the exact categories
\begin{equation} \label{lcth-W-T-lin-direct-image}
 f_!\:Y\Lcth^\lin_\bT\lrarrow X\Lcth^\lin_\bW.
\end{equation}
 We refer to~\cite[Section~3.3]{Pcosh} for the details
(see also~\cite[Section~2.4]{Pdomc}).

\subsection{The direct image-restriction adjunction}
 In the context of the proofs of the main results of the present
paper, the only morphisms of schemes that we will use are open
immersions.
 For this reason, we do not go into an (otherwise rather complicated)
discussion of inverse images of locally contraherent cosheaves with
respect to arbitrary morphisms of schemes~\cite[Sections~2.3
and~3.3]{Pcosh}, \cite[Section~2.5]{Pdomc}, \cite[Section~3.12]{Pform}.
 For open immersion morphisms, the construction of the inverse
images of locally contraherent cosheaves (as well as of quasi-coherent
sheaves) is much more straightforward.

 Let $(X,\cO_X)$ be a ringed space and $Y\subset X$ be an open
subset.
 We endow the topological space $Y$ with the sheaf of rings
$\cO_Y=\cO_X|_Y$ obtained by restricting the sheaf of rings $\cO_X$
on $X$ to the open subset $Y\subset X$.
 Then there is a natural morphism of ringed spaces
$j\:(Y,\cO_Y)\rarrow(X,\cO_X)$.

 Given a sheaf of $\cO_X$\+modules $\cM$ on $X$, the restriction
$\cM|_Y$ of the sheaf $\cM$ to the open subset $Y\subset X$ is
the sheaf of $\cO_Y$\+modules on $Y$ defined by the rule
$$
 (\cM|_Y)(V)=\cM(V)
$$
for any open subset $V\subset Y\subset X$.
 This construction defines the restriction functor
\begin{equation} \label{restriction-of-sheaves-of-O-modules}
 j^*\:(X,\cO_X)\Sh\lrarrow(Y,\cO_Y)\Sh, \qquad j^*\cM=\cM|_Y.
\end{equation}

 For any two sheaves of $\cO$\+modules $\cM$ on $X$ and $\cN$ on $Y$,
there is a natural adjunction isomorphism of abelian groups
\begin{equation} \label{sheaves-direct-image-restriction-adjunction}
 \Hom_{\cO_X}(\cM,j_*\cN)\simeq\Hom_{\cO_Y}(j^*\cM,\cN),
\end{equation}
where the notation $\Hom_{\cO_X}$ and $\Hom_{\cO_Y}$ stands for
the groups of morphisms in the categories $(X,\cO_X)\Sh$ and
$(Y,\cO_Y)\Sh$.
 In other words, the restriction functor $j^*\:(X,\cO_X)\Sh\rarrow
(Y,\cO_Y)\Sh$ \,\eqref{restriction-of-sheaves-of-O-modules} is
left adjoint to the direct image functor $j_*\:(Y,\cO_Y)\Sh\rarrow
(X,\cO_X)\Sh$ \,\eqref{direct-image-of-sheaves-of-O-modules}.

 For any scheme $X$ and any open subscheme $Y\subset X$,
the restriction functor $j^*$
\,\eqref{restriction-of-sheaves-of-O-modules} takes quasi-coherent
sheaves on $X$ to quasi-coherent sheaves on~$X$.
 So we have the restriction functor
\begin{equation} \label{restriction-of-qcoh-sheaves}
 j^*\:X\Qcoh\lrarrow Y\Qcoh, \qquad j^*\cM=\cM|_Y.
\end{equation}
 The restriction functor of quasi-coherent sheaves $j^*$
\,\eqref{restriction-of-qcoh-sheaves} is exact as a functor between
Grothendieck abelian categories.

 Dual-analogously, given a cosheaf of $\cO_X$\+modules $\fP$ on $X$,
the restriction $\fP|_Y$ of the cosheaf $\cM$ to the open subset
$Y\subset X$ is the cosheaf of $\cO_Y$\+modules on $Y$ defined by
the rule
$$
 (\fP|_Y)[V]=\fP[V]
$$
for any open subset $V\subset Y\subset X$.
 This construction defines the restriction functor
\begin{equation} \label{restriction-of-cosheaves-of-O-modules}
 j^!\:(X,\cO_X)\Cosh\lrarrow(Y,\cO_Y)\Cosh, \qquad j^!\fP=\fP|_Y.
\end{equation}

 For any two cosheaves of $\cO$\+modules $\fP$ on $X$ and $\fQ$ on $Y$,
there is a natural adjunction isomorphism of abelian groups
\begin{equation} \label{cosheaves-direct-image-restriction-adjunction}
 \Hom^{\cO_X}(j_!\fQ,\fP)\simeq\Hom^{\cO_Y}(\fQ,j^!\fP),
\end{equation}
where the notation $\Hom^{\cO_X}$ and $\Hom^{\cO_Y}$ stands for
the groups of morphisms in the categories $(X,\cO_X)\Cosh$ and
$(Y,\cO_Y)\Cosh$.
 In other words, the restriction functor $j^!\:(X,\cO_X)\Cosh\rarrow
(Y,\cO_Y)\Cosh$ \,\eqref{restriction-of-cosheaves-of-O-modules} is
right adjoint to the direct image functor $j_!\:(Y,\cO_Y)\Cosh\rarrow
(X,\cO_X)\Cosh$ \,\eqref{direct-image-of-cosheaves-of-O-modules}.
 This is~\cite[formula~(2.7) in Section~2.3]{Pcosh}
\cite[formula~(19) in Section~2.4]{Pdomc}, or~\cite[formula~(30)
in Section~3.10]{Pform}.

 Now let $X$ be a scheme with an open covering $\bW$, and let
$Y\subset X$ be an open subscheme.
 Define the restriction $\bW|_Y$ of $\bW$ to $Y$ by the rule
$\bW|_Y=\{Y\cap W\mid W\in\bW\}$.
 Clearly, $\bW|_Y$ is an open covering of~$Y$.

 Then the restriction functor $j^!\:(X,\cO_X)\Cosh\rarrow
(Y,\cO_Y)\Cosh$ \,\eqref{restriction-of-cosheaves-of-O-modules}
takes $\bW$\+locally contraherent cosheaves on $X$ to
$\bW|_Y$\+locally contraherent cosheaves on $Y$, and induces
an exact functor between the respective exact categories
\begin{equation} \label{restriction-of-lcth-cosheaves}
 j^!\:X\Lcth_\bW\lrarrow Y\Lcth_{\bW|_Y}.
\end{equation}
 Furthermore, the same restriction functor $j^!\:(X,\cO_X)\Cosh
\rarrow(Y,\cO_Y)\Cosh$ takes locally cotorsion $\bW$\+locally
contraherent cosheaves on $X$ to locally cotorsion $\bW|_Y$\+locally
contraherent cosheaves on $Y$, and induces an exact functor between
their exact categories
\begin{equation} \label{restriction-of-lcth-lct-cosheaves}
 j^!\:X\Lcth^\lct_\bW\lrarrow Y\Lcth^\lct_{\bW|_Y}.
\end{equation}
 Finally, the same functor of restriction of cosheaves of
$\cO$\+modules to the open subset takes locally injective
$\bW$\+locally contraherent cosheaves on $X$ to locally injective
$\bW|_Y$\+locally contraherent cosheaves on $Y$, and induces
an exact functor between the exact categories
\begin{equation} \label{restriction-of-lcth-lin-cosheaves}
 j^!\:X\Lcth^\lin_\bW\lrarrow Y\Lcth^\lin_{\bW|_Y}.
\end{equation}

 Notice that if, in addition, the open immersion morphism
$j\:Y\rarrow X$ is affine, then it is $(\bW,\bW|_Y)$\+affine.
 In this case, the constructions of the functors of direct image
of (locally) contraherent cosheaves from
Section~\ref{direct-image-and-contraherence-subsecn} are applicable
to~$j$.

 One typical setting in which we will apply the constructions of direct
images and restrictions of locally contraherent cosheaves in the present
paper is that of an open subscheme subordinate to the covering.
 If the open subscheme $Y\subset X$ is subordinate to the open covering
$\bW$ of $X$, then the open covering $\bW|_Y$ of $Y$ contains the whole
scheme $Y$ as one of the open subsets in the covering, i.~e.,
$Y\in\bW|_Y$.
 In this case, all the $\bW|_Y$\+locally contraherent cosheaves on $Y$
are contraherent, that is, $Y\Lcth_{\bW|_Y}=Y\Ctrh$ (and similarly for
the locally cotorsion and locally injective locally contraherent
cosheaves).


\Section{Quasi-Compact Semi-Separated Schemes}  \label{qcss-secn}

 The proof of the following proposition is based on the argument
dual-analogous to the proof of the existence of enough flat
quasi-coherent sheaves on a quasi-compact semi-separated scheme given
in~\cite[Lemma~A.1]{EP} (see also~\cite[Lemmas~4.1.1 and~4.1.8]{Pcosh}).

\begin{prop} \label{qcoh-ssep-enough-lin-prop}
 Let $X$ be a quasi-compact semi-separated scheme with an open
covering\/~$\bW$.
 Then, for any\/ $\bW$\+locally contraherent cosheaf\/ $\fP$ on $X$,
there exists a locally injective\/ $\bW$\+locally contraherent cosheaf\/
$\fJ$ on $X$ together with an admissible monomorphism\/ $\fP\rarrow\fJ$
in the exact category $X\Lcth_\bW$ of\/ $\bW$\+locally contraherent
cosheaves on~$X$.
\end{prop}

\begin{proof}
 This is a partial version of~\cite[Lemma~4.3.2(a)]{Pcosh}.
 For the sake of completeness of the exposition, we spell out
the proof below.

 Let $X=\bigcup_{\alpha=1}^N U_\alpha$ be a finite affine open covering
of $X$ subordinate to the open covering~$\bW$.
 Proceeding by induction on $\beta=0$,~\dots, $N$, we will construct
admissible monomorphisms $\fP\rarrow\fJ_\beta$ in $X\Lcth_\bW$ such
that the restriction of $\fJ_\beta$ to the open subset
$V_\beta=\bigcup_{\alpha=1}^\beta U_\alpha\subset X$ is a locally
injective $\bW|_{V_\beta}$\+locally contraherent cosheaf on~$V_\beta$.
 In the case of $\beta=0$ (the induction base), we have
$V_0=\varnothing$, so it suffices to take $\fJ_0=\fP$ with the identity
monomorphism $\fP\rarrow\fJ_0$.

 For $\beta\ge1$, assuming that an admissible monomorphism
$\fP\rarrow\fJ_{\beta-1}$ with the desired property has been
constructed already, we put $\fK=\fJ_{\beta-1}$, \ $V=V_{\beta-1}$,
and $U=U_\beta$.
 Denote the identity open immersion morphisms by $j\:U\rarrow X$
and $h\:V\rarrow X$.

 We need to use the fact that the assertion of the proposition holds
for the affine scheme $U$ with its trivial open covering $\bW_U=\{U\}$.
 Specifically, there exists an admissible monomorphism
$j^!\fK\rarrow\fI$ in $U\Ctrh$ with a locally injective contraherent
cosheaf $\fI$ on~$U$.

 Indeed, by Lemmas~\ref{ctrh-cosheaves-on-affine-scheme}
and~\ref{ctrh-lct-lin-cosheaves-on-affine-scheme}(b), we have
an equivalence of exact categories $U\Ctrh\simeq\cO_X(U)\Modl^\cta$
identifying the full subcategories $U\Ctrh^\lin\subset U\Ctrh$
and $\cO_X(U)\Modl^\inj\subset\cO_X(U)\Modl$.
 Consider the contraadjusted $\cO_X(U)$\+module $(j^!\fK)[U]=\fK[U]$,
and pick an injective $\cO_X(U)$\+module $I$ together with
an injective morphism of $\cO_X(U)$\+modules $\fK[U]\rarrow I$.
 Then the quotient $\cO_X(U)$\+module $I/\fK[U]$ is contraadjusted,
so $\fK[U]\rarrow I$ is an admissible monomorphism in
$\cO_X(U)\Modl^\cta$.
 Let $j^!\fK\rarrow\fI$ be the related admissible monomorphism in
$U\Ctrh$.

 So we have an (admissible) short exact sequence $0\rarrow j^!\fK
\rarrow\fI\rarrow\fR\rarrow0$ in $U\Ctrh$.
 Applying the exact functor of direct image of contraherent
cosheaves $j_!$ \,\eqref{ctrh-direct-image} for the affine open
immersion morphism $j\:U\rarrow X$ (here we use the assumption that
the scheme $X$ is semi-separated), we obtain a short exact sequence
$0\rarrow j_!j^!\fK\rarrow j_!\fI\rarrow j_!\fR\rarrow0$ in
$X\Ctrh\subset X\Lcth_\bW$.
 The adjunction isomorphism of abelian
groups~\eqref{cosheaves-direct-image-restriction-adjunction} provides
an adjunction morphism $j_!j^!\fK\rarrow\fK$ in $X\Lcth_\bW$.
 Taking the pushout of this short exact sequence with respect to
this morphism, we obtain a short exact sequence
$0\rarrow\fK\rarrow\fL\rarrow j_!\fR\rarrow0$ in $X\Lcth_\bW$.

 Put $\fJ_\beta=\fL$.
 We claim that the $\bW$\+locally contraherent cosheaf $\fL$ on $X$
is locally injective in restriction to the open subscheme
$V_\beta=V\cup U$.
 In view of the discussion in Section~\ref{lct-lin-subsecn}, it
suffices to check that the contraherent cosheaf $j^!\fL$ is
locally injective on $U$ and the $\bW|_V$\+locally contraherent
cosheaf $h^!\fL$ is localy injective on~$V$.

 The restriction functor~$j^!$ takes the adjunction morphism
$j_!j^!\fK\rarrow\fK$ to an isomorphism.
 Hence we have $j^!\fL\simeq j^!j_!\fI\simeq\fI$, and $\fI$ is
a locally injective contraherent cosheaf on $U$ by construction.
 This takes care of the restriction to~$U$.

 The proof of the local injectivity of the locally contraherent
cosheaf $h^!\fL$ on $V$ is a bit more involved.
 Consider the intersection $U\cap V\subset X$, and denote
the identity open immersion morphisms by $j'\:U\cap V\rarrow V$
and $h'\:U\cap V\rarrow U$.

 We have a base change isomorphism $h^!j_!\fR\simeq j'_!h'{}^!\fR$,
which is obvious immediately from the construction of the direct image.
 Restricting the short exact sequence $0\rarrow j^!\fK\rarrow\fI
\rarrow\fR\rarrow0$ in $U\Ctrh$ to the open subscheme
$U\cap V\subset U$, we obtain a short exact sequence
$0\rarrow h'{}^!j^!\fK\rarrow h'{}^!\fI\rarrow h'{}^!\fR\rarrow0$
of contraherent cosheaves on $U\cap V$ (see
formula~\eqref{restriction-of-lcth-cosheaves}).

 Now the $\bW|_V$\+locally contraherent cosheaf $h^!\fK$ on $V$
is locally injective by the induction assumption; hence
the contraherent cosheaf $h'{}^!j^!\fK\simeq j'{}^!h^!\fK$ on
$U\cap V$ is locally injective, too.
 The contraherent cosheaf $h'{}^!\fI$ on $U\cap V$ is also locally
injective as a restriction of the locally injective contraherent
cosheaf $\fI$ to an open subscheme.
 Since the full subcategory $(U\cap\nobreak V)\Ctrh^\lin$ is
closed under cokernels of admissible monomorphisms in
$(U\cap\nobreak V)\Ctrh$ (see Section~\ref{lct-lin-subsecn}),
it follows that the contraherent cosheaf $h'{}^!\fR$ on $U\cap V$
is locally injective as well.
 Since the direct image functor~$j'_!$ with respect to the flat affine
morphism $j'\:U\cap V\rarrow V$ preserves the local injectivity
of contraherent cosheaves by formula~\eqref{ctrh-lin-direct-image},
we arrive to the conclusion that the contraherent cosheaf
$h^!j_!\fR\simeq j'_!h'{}^!\fR$ on $V$ is locally injective.

 Finally, restricting to the open subscheme $V\subset X$ the short
exact sequence $0\rarrow\fK\rarrow\fL\rarrow j_!\fR\rarrow0$ in
$X\Lcth_\bW$, we obtain a short exact sequence $0\rarrow h^!\fK\rarrow
h^!\fL\rarrow h^!j_!\fR\rarrow0$ in $V\Lcth_{\bW|_V}$ with locally
injective $\bW|_V$\+locally contraherent cosheaves $h^!\fK$ and
$h^!j_!\fR$.
 As the full subcategory $V\Lcth^\lin_{\bW|_V}$ is closed under
extensions in $V\Lcth_{\bW|_V}$ (see Section~\ref{lct-lin-subsecn}),
it follows that the $\bW|_V$\+locally contraherent cosheaf $h^!\fL$
on $V$ is also locally injective, as desired.
\end{proof}

\begin{rem}
 For any quasi-compact semi-separated scheme $X$ with an open covering
$\bW$, the full subcategory of locally injective $\bW$\+locally
contraherent cosheaves $X\Lcth^\lin_\bW$ is, actually, the right-hand
class of a hereditary complete cotorsion pair in the exact category
$X\Lcth_\bW$.
 The objects from the left-hand class are called the \emph{antilocal}
contraherent cosheaves on $X$ (all such objects are globally
contraherent on $X$, and the full subcategory of antilocal contraherent
cosheaves on $X$ does not depend on the open covering~$\bW$).
 This is explained in~\cite[Section~4.3]{Pcosh} (see
also~\cite[Sections~4.4\+-4.5]{Pdomc}.
\end{rem}


\Section{Main Negative Result}

 We start with a fairly standard lemma about quasi-coherent sheaves.

\begin{lem} \label{quasi-compact-intersected-with-affine}
 Let $X$ be an affine scheme, $V\subset X$ be a quasi-compact open
subscheme, and $U\subset X$ be an affine open subscheme.
 Let\/ $\cM$ be a quasi-coherent sheaf on~$X$.
 Then the restriction map of\/ $\cO(X)$\+modules\/
$\cM(V)\rarrow\cM(U\cap V)$ induces an isomorphism of\/
$\cO_X(U)$\+modules
$$
 \cO_X(U)\ot_{\cO(X)}\cM(V)\simeq\cM(U\cap V).
$$
\end{lem}

\begin{proof}
 Let $V=\bigcup_{\alpha=1}^N V_\alpha$ be a finite affine open covering
of the scheme~$V$.
 Then the sheaf axiom~\eqref{sheaf-axiom-for-topology-base} tells us
that the $\cO(X)$\+module $\cM(V)$ can be computed as the kernel of
the $\cO(X)$\+module map
\begin{equation} \label{M-V-sheaf-axiom-computing}
 \bigoplus\nolimits_{\alpha=1}^N \cM(V_\alpha)\lrarrow
 \bigoplus\nolimits_{1\le\alpha<\beta\le N} \cM(V_\alpha\cap V_\beta).
\end{equation}
 Now $U\cap V=\bigcup_{\alpha=1}^N U\cap V_\alpha$ is a finite affine
open covering of the scheme $U\cap V$.
 Hence the $\cO_X(U)$\+module $\cM(U\cap V)$ can be similarly computed
as the kernel of the $\cO_X(U)$\+module map
\begin{equation} \label{M-U-cap-V-sheaf-axiom-computing}
 \bigoplus\nolimits_{\alpha=1}^N \cM(U\cap V_\alpha)\lrarrow
 \bigoplus\nolimits_{1\le\alpha<\beta\le N}
 \cM(U\cap V_\alpha\cap V_\beta).
\end{equation}
 It remains to observe that
the map~\eqref{M-U-cap-V-sheaf-axiom-computing} can be obtained by
applying the functor $\cO_X(U)\ot_{\cO(X)}{-}$ to
the map~\eqref{M-V-sheaf-axiom-computing}, and that this functor
preserves kernels (since the $\cO(X)$\+module $\cO_X(U)$ is flat,
and in fact, very flat).
\end{proof}

\begin{lem} \label{co-generation-square-lemma}
 Let $R\rarrow S$ be a homomorphism of associative rings. \par
\textup{(a)} Suppose given a commutative square diagram of
left $R$\+module maps
\begin{equation} \label{generation-square-diagram}
\begin{gathered}
 \xymatrix{
  F \ar[r]^-q \ar@{->>}[d]_f & F' \ar[d]^{f'} \\
  M \ar[r]_-r & S\ot_RM 
 }
\end{gathered}
\end{equation}
where the $R$\+module map\/ $r\:M\rarrow S\ot_RM$ is induced
by the ring homomorphism $R\rarrow S$, the $R$\+module map
$f\:F\rarrow M$ is surjective, $F'$ is the underlying $R$\+module of
an $S$\+module, and\/ $f'\:F'\rarrow S\ot_RM$ is an $S$\+module map.
 Then the map $f'\:F'\rarrow S\ot_RM$ is surjective. \par
\textup{(b)} Suppose given a commutative square diagram of
left $R$\+module maps
\begin{equation} \label{cogeneration-square-diagram}
\begin{gathered}
 \xymatrix{
  J & J' \ar[l]_-d \\
  G \ar@{>->}[u]^g & \Hom_R(S,G) \ar[l]^-c \ar[u]_{g'}
 }
\end{gathered}
\end{equation}
where the $R$\+module map\/ $c\:\Hom_R(S,G)\rarrow G$ is induced
by the ring homomorphism $R\rarrow S$, the $R$\+module map
$g\:G\rarrow J$ is injective, $J'$ is the underlying $R$\+module of
an $S$\+module, and\/ $g'\:\Hom_R(S,G)\rarrow J'$ is an $S$\+module map.
 Then the map $g'\:\Hom_R(S,G)\rarrow J'$ is injective.
\end{lem}

\begin{proof}
 Part~(a), which is included here as an illustration for part~(b),
is clear and intuitive.
 The point is that the $S$\+module $S\ot_RM$ is generated by
the image of the map~$r$, which is contained in the image of
the $S$\+module map~$f'$.
 Let us prove part~(b), which we will really use.

 Let $v\:S\rarrow G$ be a left $R$\+module map; so $v\in\Hom_R(S,G)$.
 Assume that $g'(v)=0$.
 Then we also have $g'(sv)=0$ for all $s\in S$, since $g'$~is
an $S$\+module map.
 Since the map~$g$ is injective, it follows that $c(sv)=0$.
 By construction, we have $c(sv)=(sv)(1)=v(s)\in G$.
 Thus $v(s)=0$ for all $s\in S$; so $v=0$.
\end{proof}

 The following lemma is a dual analogue of~\cite[Lemma~2.1]{SS}.

\begin{lem} \label{slavik-stovicek-lemma}
 Let $X$ be an affine scheme, $U\subset X$ be a quasi-compact open
subscheme, and\/ $\fJ\in X\Ctrh^\lin$ be a locally injective
contraherent cosheaf on~$X$.
 Consider the homomorphism of\/ $\cO_X(U)$\+modules
\begin{equation} \label{qcomp-in-affine-ctrh-module-map}
 \fJ[U]\lrarrow\Hom_{\cO(X)}(\cO_X(U),\fJ[X])
\end{equation}
induced by the corestriction map of\/ $\cO(X)$\+modules\/
$\fJ[U]\rarrow\fJ[X]$.
 Then the map~\eqref{qcomp-in-affine-ctrh-module-map} is
an isomorphism.
\end{lem}

\begin{proof}
 Put $R=\cO(X)$.
 By Lemma~\ref{ctrh-lct-lin-cosheaves-on-affine-scheme}(b),
locally injective contraherent cosheaves $\fJ$ on $X$ correspond to
injective $R$\+modules $J=\fJ[X]$.
 For any injective $R$\+module $J$, there exists a flat (and even
free/projective) $R$\+module $F$ such that $J$ is a direct summand of
the $R$\+module $F^+=\Hom_\boZ(F,\boQ/\boZ)$.
 Denoting by $\cF$ the flat/projective/free quasi-coherent sheaf on
$X$ corresponding to the $R$\+module $F$, we conclude that $\fJ$ is
a direct summand of the locally injective contraherent cosheaf $\cF^+$
on $X$, in the notation of
Remarks~\ref{finite-coverings-suffice-remarks}(b\+-d)
and~\ref{character-cosheaf-ctrh-lct-remark}.
 So it suffices to consider the case of $\fJ=\cF^+$.
 In this case, the desired
isomorphism~\eqref{qcomp-in-affine-ctrh-module-map} can be obtained
by applying the character module functor
$({-})^+=\Hom_\boZ({-},\boQ/\boZ)$ to the isomorphism of
$\cO_X(U)$\+modules $\cO_X(U)\ot_{\cO(X)}\cF(X)\simeq\cF(U)$
from~\cite[Lemma~2.1]{SS}.
\end{proof}

 The following theorem is the main result of this paper.
 It is a dual analogue of the theorem of Sl\'avik
and \v St\!'ov\'\i\v cek~\cite[Theorem~2.2]{SS}.

\begin{thm} \label{non-ssep-not-enough-lin-theorem}
 Let $X$ be a quasi-compact, quasi-separated scheme that is \emph{not}
semi-separated.
 Then there is a locally cotorsion contraherent cosheaf\/
$\fG\in X\Ctrh^\lct$ on $X$ for which there \emph{does not exist}
an admissible monomorphism\/ $\fG\rarrow\fJ$ in the exact category
$X\Lcth$ or $X\Lcth^\lct$ from\/ $\fG$ into any locally injective
locally contraherent cosheaf\/ $\fJ\in X\Lcth^\lin$.
\end{thm}

\begin{proof}
 We follow the argument from~\cite[proof of Theorem~2.2]{SS}, with
suitable enhancements.
 Let $U$, $V\subset X$ be two affine open subschemes whose intersection
$T=U\cap V$ is not affine.
 Since $X$ is quasi-separated, the scheme $T$ is quasi-compact.
 Put $R=\cO_X(U)$, so $U=\Spec R$.
 Then there is a finite collection of elements $f_1$,~\dots, $f_n\in R$
such that $T=\Spec R[f_1^{-1}]\cup\dotsb\cup\Spec R[f_n^{-1}]\subset U$.

 Let $I=(f_1,\dotsc,f_n)\subset R$ be the ideal spanned by $f_1$,
\dots, $f_n$ in $R$, and let $\cI$ be the quasi-coherent sheaf
(of ideals) on $U$ corresponding to the $R$\+module~$I$.
 Denote the identity open immersion morphism by $j\:U\rarrow X$,
and consider the locally cotorsion contraherent cosheaf
$\fG=(j_*\cI)^+\simeq j_!(\cI^+)$ on~$X$ (in the notation of
Remark~\ref{character-cosheaf-ctrh-lct-remark}).
 Assuming that there exists an open covering $\bW$ of $X$,
a locally injective $\bW$\+locally contraherent cosheaf $\fJ$ on $X$,
and an admissible monomorphism $\fG\rarrow\fJ$ in $X\Lcth_\bW$,
we will come to a contradiction.

 Let $U=\bigcup_\alpha U_\alpha$ be a covering of $U$ by affine open
subschemes $U_\alpha$ subordinate to $\bW$, and let $V=\bigcup_\beta
V_\beta$ be a covering of $V$ by affine open subschemes $V_\beta$
subordinate to~$\bW$.
 Our next aim is to prove that there exists a pair of indices
$(\alpha,\beta)$ such that the open subscheme $U_\alpha\cap V_\beta
\subset X$ is not affine.

 We will consider Cartesian products of schemes over $\Spec\boZ$;
so, throughout the rest of this proof, the Cartesian product
sign~$\times$ without subindex means $\times_{\Spec\boZ}$.
 The pair of open immersion morphisms $T\rarrow U$ and $T\rarrow V$
induces a morphism of schemes $\delta\:T\rarrow U\times V$.
 The affine scheme $U\times V$ is covered by its affine open
subschemes $U_\alpha\times V_\beta$, that is, $U\times V=
\bigcup_{\alpha,\beta}(U_\alpha\times V_\beta)$.
 For any pair of indices $\alpha$ and~$\beta$, we have
$\delta^{-1}(U_\alpha\times V_\beta)=U_\alpha\cap V_\beta\subset T$.
 Since the scheme $T$ is not affine, the scheme morphism~$\delta$
cannot be affine.
 By~\cite[D\'efinition~II.1.2.1 and Corollaire~II.1.3.2]{EGAII}
or~\cite[Lemma Tag~01S8]{SP}, it follows that a desired pair
of indices $\alpha$ and~$\beta$ with a nonaffine intersection
$U_\alpha\cap V_\beta$ exists.

 For the chosen pair of indices $\alpha$ and~$\beta$, we put
$U'=U_\alpha$, \ $V'=V_\beta$, and $T'=U_\alpha\cap V_\beta$.
 Then the open immersion morphism $T'\rarrow T$ is affine as
a base change of the morphism of affine schemes $U'\times V'
\rarrow U\times V$; specifically,
$T'=(U'\times V')\times_{(U\times V)}T$.
 Hence, restricting to $T'$ the affine open covering
$T=\bigcup_{i=1}^n\Spec R[f_i^{-1}]$, we obtain
an affine open covering $T'=\bigcup_{i=1}^n(T'\cap\Spec R[f_i^{-1}])$
of the scheme~$T'$.
 To repeat, the point is that the open subschemes
$T'\cap\Spec R[f_i^{-1}]$ are affine.
 Hence, according to~\cite[Chapter~II, Section~II.2,
Exercise~2.17(b)]{HarAG}, the fact that the scheme $T'$ is not affine
implies that the restrictions of the elements $f_1$,~\dots,
$f_n\in\cO_X(U)$ to the open subscheme $T'\subset U$ do \emph{not}
generate the unit ideal of the ring $\cO_X(T')$.

 The following commutative square diagram is a dual analogue
of the leftmost square of the main commutative diagram
in~\cite[proof of Theorem~2.2]{SS}:
\begin{equation} \label{slavik-stovicek-leftmost-square}
\begin{gathered}
 \xymatrix{
  \fJ[V'] & \fJ[U\cap V'] \ar[l] \\
  \fG[V'] \ar@{>->}[u] \ar@{=}[r] & \fG[U\cap V'] \ar[u]
 }
\end{gathered}
\end{equation}
 Here the upper horizontal map $\fJ[U\cap V']\rarrow\fJ[V']$
is the corestriction map in the locally injective locally
contraherent cosheaf~$\fJ$.
 The corestriction map $\fG[U\cap V']\rarrow\fG[V']$ in
the contraherent cosheaf $\fG=(j_*\cI)^+$ is an isomorphism because
the restriction map $(j_*\cI)(V')\rarrow(j_*\cI)(V'\cap U)$ is
an isomorphism by the construction of the direct image functor
$j_*\:U\Qcoh\rarrow X\Qcoh$.
 We are using the fact that the scheme $U\cap V'$ is quasi-compact
and semi-separated here.

 The map of cosections $\fG[V']\rarrow\fJ[V']$ induced by
the admissible monomorphism of $\bW$\+locally contraherent cosheaves
$\fG\rarrow\fJ$ on $X$ is injective because $V'\subset X$ is
an affine open subscheme subordinate to~$\bW$.
 It follows from the commutativity of
the diagram~\eqref{slavik-stovicek-leftmost-square} that the map of
cosections $\fG[U\cap V']\rarrow\fJ[U\cap V']$ induced by the morphism
of cosheaves $\fG\rarrow\fJ$ is injective, too.

 Now let us consider the commutative square diagram
\begin{equation} \label{inserted-middle-square}
\begin{gathered}
 \xymatrix{
  \fJ[U\cap V'] & \fJ[U'\cap V'] \ar[l] \\
  \fG[U\cap V'] \ar@{>->}[u] & \fG[U'\cap V'] \ar[l] \ar[u]
 }
\end{gathered}
\end{equation}
 Here the horizontal maps are the corestriction maps in the cosheaves
$\fJ$ and $\fG$, while the vertical maps are the maps of cosections
induced by the morphism of cosheaves $\fG\rarrow\fJ$.
 The leftmost vertical map in~\eqref{inserted-middle-square} is
injective according to the argument above.

 Put $\cM=j_*\cI$; so $\fG=\cM^+$.
 The corestriction map $\fG[U'\cap V']\rarrow\fG[U\cap V']$ is
obtained by applying the character module functor $({-})^+$
to the restriction map $\cM(U\cap V')\rarrow\cM(U'\cap V')$.
 Now we use Lemma~\ref{quasi-compact-intersected-with-affine} for
the affine scheme $U$, its quasi-compact open subscheme $U\cap V'
\subset U$, and the affine open subscheme $U'\subset U$.
 According to the lemma, we have an isomorphism of $\cO_X(U')$\+modules
$\cO_X(U')\ot_{\cO_X(U)}\cM(U\cap V')\simeq\cM(U'\cap V')$.
 Applying the functor $({-})^+$, we obtain an isomorphism
$$
 \fG[U'\cap V']\simeq\Hom_{\cO_X(U)}(\cO_X(U'),\>\fG[U\cap V']).
$$

 Put $R=\cO_X(U)$ and $S=\cO_X(U')$.
 Now the diagram~\eqref{inserted-middle-square} is a special case
of the diagram~\eqref{cogeneration-square-diagram}
from Lemma~\ref{co-generation-square-lemma}(b).
 According to that lemma, it follows that the rightmost vertical
map $\fG[U'\cap V']\rarrow\fJ[U'\cap V']$
in~\eqref{inserted-middle-square} is injective as well.

 Finally, the following commutative square diagram is a dual analogue
of the rightmost square of the main commutative diagram
in~\cite[proof of Theorem~2.2]{SS}:
\begin{equation} \label{slavik-stovicek-rightmost-square}
\begin{gathered}
 \xymatrix{
  \fJ[U'\cap V'] \ar[r] & \fJ[U'] \\
  \fG[U'\cap V'] \ar@{>->}[u] \ar[r] & \fG[U'] \ar@{>->}[u]
 }
\end{gathered}
\end{equation}
 Once again, the horizontal maps are the corestriction maps in
the cosheaves $\fJ$ and $\fG$, while the vertical maps are the maps of
cosections induced by the morphism of cosheaves $\fG\rarrow\fJ$.
 The leftmost vertical map
in~\eqref{slavik-stovicek-rightmost-square} is injective according to
the argument above, while the rightmost vertical map
in~\eqref{slavik-stovicek-rightmost-square} is injective because
$\fG\rarrow\fJ$ is an admissible monomorphism in $X\Lcth_\bW$ and
$U'\subset X$ is an affine open subscheme subordinate to~$\bW$.

 The argument finishes similarly to the proof in~\cite{SS}.
 The commutative diagram of $\cO_X(U')$\+module
maps~\eqref{slavik-stovicek-rightmost-square} induces a commutative
diagram of $\cO_X(U'\cap\nobreak V')$\+module maps
\begin{equation} \label{slavik-stovicek-subsequent-diagram}
\begin{gathered}
 \xymatrix{
  \fJ[U'\cap V'] \ar[r]
   & \Hom_{\cO_X(U')}(\cO_X(U'\cap V'),\>\fJ[U']) \\
  \fG[U'\cap V'] \ar@{>->}[u] \ar[r]
   & \Hom_{\cO_X(U')}(\cO_X(U'\cap V'),\>\fG[U']) \ar@{>->}[u]
 }
\end{gathered}
\end{equation}
 The upper horizontal map in~\eqref{slavik-stovicek-subsequent-diagram}
is an isomorphism by Lemma~\ref{slavik-stovicek-lemma} applied to
the quasi-compact open subscheme $U'\cap V'$ in the affine scheme $U'$
and the locally injective contraherent cosheaf $\fJ|_{U'}$ on~$U'$.
 Since the leftmost vertical map
in~\eqref{slavik-stovicek-subsequent-diagram} is injective, it follows
that the lower horizontal map is injective, too.

 The lower horizontal map $\fG[U'\cap V']\rarrow
\Hom_{\cO_X(U')}(\cO_X(U'\cap V'),\>\fG[U'])$
in~\eqref{slavik-stovicek-subsequent-diagram} can be obtained by
applying the functor $({-})^+$ to the natural map
\begin{equation} \label{slavik-stovicek-surjective-map}
 \cO_X(U'\cap V')\ot_{\cO_X(U')}\cI(U')\lrarrow\cI(U'\cap V').
\end{equation}
 Consequently, the map~\eqref{slavik-stovicek-surjective-map} is
surjective.

 Now we recall that $\cI$ is a quasi-coherent sheaf of ideals on
$U=\Spec R$.
 By construction, the restriction of $\cI$ to $\Spec R[f_i^{-1}]
\subset U$ coincides with the structure sheaf of $\Spec R[f_i^{-1}]$
for every $1\le i\le n$.
 Therefore, the restriction of the sheaf of ideals $\cI$ to
the open subscheme $U\cap V=\bigcup_{i=1}^n\Spec R[f_i^{-1}]$
coincides with the structure sheaf of $U\cap V$.
 So we have $\cI(U'\cap V')=\cO_X(U'\cap V')$.

 On the other hand, we have $\cI(U')=\cO_X(U')\ot_{\cO_X(U)}\cI(U)$.
 Hence the map~\eqref{slavik-stovicek-surjective-map} is isomorphic
to the natural map
\begin{equation} \label{surjective-map-rewritten}
 \cO_X(U'\cap V')\ot_{\cO_X(U)}\cI(U)\lrarrow\cO_X(U'\cap V').
\end{equation}
 Recall also the notation $\cI(U)=I=(f_1,\dotsc,f_n)\subset R$
and $T'=U'\cap V'$.
 It remains to point out that surjectivity of
the map~\eqref{surjective-map-rewritten} contradicts our previous
observation that restrictions of the elements $f_1$,~\dots,
$f_n\in\cO_X(U)$ to the open subscheme $T'\subset U$ do not
generate the unit ideal of the ring $\cO_X(T')$.
 We have arrived to a contradiction, proving the theorem.
\end{proof}

 The following corollary summarizes the results of
Section~\ref{qcss-secn} and the present section, together
with~\cite[Theorem~2.2]{SS}.

\begin{cor} \label{slavik-stovicek-style-main-corollary}
 Let $X$ be a quasi-compact, quasi-separated scheme with an open
covering\/~$\bW$.
 Then the following conditions are equivalent:
\begin{enumerate}
\item the scheme $X$ is semi-separated;
\item every quasi-coherent sheaf on $X$ is a quotient sheaf of
a flat quasi-coherent sheaf;
\item for every contraherent cosheaf\/ $\fP\in X\Ctrh$, there exists
an admissible monomorphism\/ $\fP\rarrow\fJ$ in the exact category
$X\Ctrh$ from\/ $\fP$ to a locally injective contraherent cosheaf\/
$\fJ\in X\Ctrh^\lin$;
\item for every locally cotorsion contraherent cosheaf\/
$\fP\in X\Ctrh^\lct$, there exists an admissible monomorphism\/
$\fP\rarrow\fJ$ in the exact category $X\Ctrh^\lct$ (or equivalently,
in $X\Ctrh$) from\/ $\fP$ to a locally injective contraherent
cosheaf\/ $\fJ\in X\Ctrh^\lin$;
\item for every\/ $\bW$\+locally contraherent cosheaf\/
$\fP\in X\Lcth_\bW$, there exists an admissible monomorphism\/
$\fP\rarrow\fJ$ in the exact category $X\Lcth_\bW$ from\/ $\fP$
to a locally injective\/ $\bW$\+locally contraherent cosheaf\/
$\fJ\in X\Lcth^\lin_\bW$;
\item for every locally cotorsion\/ $\bW$\+locally contraherent
cosheaf\/ $\fP\in X\Lcth^\lct_\bW$, there exists an admissible
monomorphism\/ $\fP\rarrow\fJ$ in the exact category $X\Lcth^\lct_\bW$
(or equivalently, in $X\Lcth_\bW$) from\/ $\fP$ to a locally injective\/
$\bW$\+locally contraherent cosheaf\/ $\fJ\in X\Lcth^\lin_\bW$;
\item for every contraherent cosheaf\/ $\fP\in X\Ctrh$, there exists
an admissible monomorphism\/ $\fP\rarrow\fJ$ in the exact category
$X\Lcth$ from\/ $\fP$ to a locally injective locally contraherent
cosheaf\/ $\fJ\in X\Lcth^\lin$;
\item for every locally cotorsion contraherent cosheaf\/
$\fP\in X\Ctrh^\lct$, there exists an admissible monomorphism\/
$\fP\rarrow\fJ$ in the exact category $X\Lcth^\lct$ (or equivalently,
in $X\Lcth$) from\/ $\fP$ to a locally injective locally contraherent
cosheaf\/ $\fJ\in X\Lcth^\lin$.
\end{enumerate}
\end{cor}

\begin{proof}
 (1)~$\Longrightarrow$~(2) This is~\cite[Section~2.4]{M-n},
\cite[Section~3.2]{M-th}, \cite[Lemma~A.1]{EP},
or~\cite[Lemma~4.1.8 or Corollary~4.1.11(a)]{Pcosh}
(see~\cite[Lemma~4.1.1 or Corollary~4.1.4(a)]{Pcosh} for
a stronger very flat version).

 (2)~$\Longrightarrow$~(1) This is~\cite[Theorem~2.2]{SS}.

 (3)~$\Longrightarrow$~(4) Holds because the full exact subcategory
$X\Ctrh^\lct$ is closed under cokernels of admissible monomorphisms
in $X\Ctrh$.

 (5)~$\Longrightarrow$~(6) Holds because the full exact subcategory
$X\Lcth^\lct_\bW$ is closed under cokernels of admissible monomorphisms
in $X\Lcth_\bW$.

 (7)~$\Longrightarrow$~(8) Holds because the full exact subcategory
$X\Lcth^\lct$ is closed under cokernels of admissible monomorphisms
in $X\Lcth$.

 (5)~$\Longrightarrow$~(3) Take $\bW=\{X\}$.
 
 (6)~$\Longrightarrow$~(4) Take $\bW=\{X\}$.
 
 (3)~$\Longrightarrow$~(7) Holds due to the inclusion of exact
categories $X\Ctrh\subset X\Lcth$.

 (4)~$\Longrightarrow$~(8) Holds due to the inclusion of exact
categories $X\Ctrh^\lct\subset X\Lcth^\lct$.

 (1)~$\Longrightarrow$~(5) This is
Proposition~\ref{qcoh-ssep-enough-lin-prop}.

 (8)~$\Longrightarrow$~(1) This is
Theorem~\ref{non-ssep-not-enough-lin-theorem}.
\end{proof}


\bigskip
\begin{thebibliography}{99}
\smallskip

\bibitem{BBE}
 L.~Bican, R.~El Bashir, E.~Enochs.
   All modules have flat covers.
\textit{Bulletin of the London Math.\ Society} \textbf{33}, \#4,
p.~385--390, 2001.

\bibitem{Bueh}
 T.~B\"uhler.
   Exact categories.
\textit{Expositiones Math.}\ \textbf{28}, \#1, p.~1--69, 2010.
\texttt{arXiv:0811.1480 [math.HO]}

\bibitem{EP}
 A.~I.~Efimov, L.~Positselski.
   Coherent analogues of matrix factorizations and relative
singularity categories.
\textit{Algebra and Number Theory} \textbf{9}, \#5, p.~1159--1292,
2015.  \texttt{arXiv:1102.0261 [math.CT]}

\bibitem{ET}
 P.~C.~Eklof, J.~Trlifaj.
   How to make Ext vanish.
\textit{Bull.\ of the London Math.\ Soc.}\ \textbf{33}, \#1,
p.~41--51, 2001.

\bibitem{En2}
 E.~Enochs.
   Flat covers and flat cotorsion modules.
\textit{Proc.\ of the Amer.\ Math.\ Soc.}\ \textbf{92}, \#2,
p.~179--184, 1984.

\bibitem{EE}
 E.~Enochs, S.~Estrada.
   Relative homological algebra in the category of quasi-coherent
sheaves.
\textit{Advances in Math.}\ \textbf{194}, \#2, p.~284--295, 2005.

\bibitem{GT}
 R.~G\"obel, J.~Trlifaj.
   Approximations and endomorphism algebras of modules.
Second Revised and Extended Edition.
De Gruyter Expositions in Mathematics 41,
De Gruyter, Berlin--Boston, 2012.

\bibitem{EGAII}
 A.~Grothendieck, J.~Dieudonn\'e.
   \'El\'ements de g\'eom\'etrie alg\'ebrique:~II.
\'Etude globale \'el\'ementaire de quelques classes de morphismes.
\textit{Publications Mathematiques de l'IH\'ES} \textbf{8},
p.~5--222, 1961.

\bibitem{EGAI}
 A.~Grothendieck, J.~A.~Dieudonn\'e.
   \'El\'ements de g\'eom\'etrie alg\'ebrique.~I.
Grundlehren der mathematischen Wissenschaften, 166.
Springer-Verlag, Berlin--Heidelberg--New York, 1971.

\bibitem{HarAG}
 R.~Hartshorne.
   Algebraic geometry.
Graduate Texts in Math., 52, Springer-Verlag,
New York--Heidelberg, 1977.

\bibitem{SP}
 A.~J.~de~Jong et al.
   The Stacks Project.
Available from \texttt{https://stacks.math.columbia.edu/}

\bibitem{Kan}
 R.~Kanda.
   Non-exactness of direct products of quasi-coherent sheaves.
\textit{Documenta Math.}\ \textbf{24}, p.~2037--2056, 2019.
\texttt{arXiv:1810.08752 [math.CT]}

\bibitem{Kra}
 H.~Krause.
   The stable derived category of a Noetherian scheme.
\textit{Compositio Math.}\ \textbf{141}, \#5, p.~1128--1162, 2005.
\texttt{arXiv:math.AG/0403526}

\bibitem{M-n}
 D.~Murfet.
   Derived categories of quasi-coherent sheaves.
Notes, October~2006.
Available from \texttt{http://www.therisingsea.org/notes}

\bibitem{M-th}
 D.~Murfet.
   The mock homotopy category of projectives and Grothendieck duality.
Ph.~D.\ Thesis, Australian National University, September 2007.
Available from \texttt{http://www. therisingsea.org/thesis.pdf}
{\hbadness=1900\par}

\bibitem{Pcosh}
 L.~Positselski.
   Contraherent cosheaves on schemes.
Electronic preprint \texttt{arXiv:1209.2995v24 [math.CT]}.

\bibitem{Pcta}
 L.~Positselski.
   Contraadjusted modules, contramodules, and reduced cotorsion modules.
\textit{Moscow Math.\ Journ.}\ \textbf{17}, \#3, p.~385--455, 2017.
\texttt{arXiv:1605.03934 [math.CT]}

\bibitem{Pctrl}
 L.~Positselski.
   An explicit self-dual construction of complete cotorsion pairs
in the relative context.
\textit{Rendiconti Semin.\ Matem.\ Univ.\ Padova} \textbf{149},
p.~191--253, 2023.  \texttt{arXiv:2006.01778 [math.RA]}

\bibitem{Pphil}
 L.~Positselski.
   Philosophy of contraherent cosheaves.
Electronic preprint \texttt{arXiv:2311.14179 [math.AG]}.

\bibitem{Pal}
 L.~Positselski.
   Local, colocal, and antilocal properties of modules and complexes
over commutative rings.
\textit{Journ.\ of Algebra} \textbf{646}, p.~100--155, 2024.
\texttt{arXiv:2212.10163 [math.AC]}

\bibitem{Pdomc}
 L.~Positselski.
   $\mathcal D$\+$\Omega$ duality on the contra side.
Electronic preprint \texttt{arXiv:2504.18460v12 [math.AG]}.

\bibitem{Pform}
 L.~Positselski.
   Contraherent cosheaves of contramodules on Noetherian formal schemes.
Electronic preprint \texttt{arXiv:2603.27732v5 [math.AG]}.

\bibitem{PSl1}
 L.~Positselski, A.~Sl\'avik.
   Flat morphisms of finite presentation are very flat.
\textit{Annali di Matem.\ Pura ed Appl.}\ \textbf{199}, \#3,
p.~875--924, 2020.  \texttt{arXiv:1708.00846 [math.AC]}

\bibitem{Sal}
 L.~Salce.
   Cotorsion theories for abelian groups.
\textit{Symposia Math.}\ \textbf{XXIII},
Academic Press, London--New York, 1979, p.~11--32.

\bibitem{SaoSt}
 M.~Saor\'\i n, J.~\v St\!'ov\'\i\v cek.
   On exact categories and applications to triangulated adjoints
and model structures.
\textit{Advances in Math.}\ \textbf{228}, \#2, p.~968--1007, 2011.
\texttt{arXiv:1005.3248 [math.CT]}

\bibitem{SS}
 A.~Sl\'avik, J.~\v St\!'ov\'\i\v cek.
   On flat generators and Matlis duality for quasicoherent sheaves.
\textit{Bulletin of the London Math.\ Soc.}\ \textbf{53}, \#1,
p.~63--74, 2021.  \texttt{arXiv:1902.05740 [math.AG]}

\bibitem{ST}
 A.~Sl\'avik, J.~Trlifaj.
   Very flat, locally very flat, and contraadjusted modules.
\textit{Journ.\ of Pure and Appl.\ Algebra} \textbf{220}, \#12,
p.~3910--3926, 2016.  \texttt{arXiv:1601.00783 [math.AC]}

\bibitem{Sto-ICRA}
 J.~\v St\!'ov\'\i\v cek.
   Exact model categories, approximation theory, and cohomology of
quasi-coherent sheaves.
\textit{Advances in representation theory of algebras}, p.~297--367,
EMS Ser.\ Congr.\ Rep., Eur.\ Math.\ Soc., Z\" urich, 2013.
\texttt{arXiv:1301.5206 [math.CT]}

\bibitem{Sto}
 J.~\v St\!'ov\'\i\v cek.
   Derived equivalences induced by big cotilting modules. 
\textit{Advances in Math.}\ \textbf{263}, p.~45--87, 2014.
\texttt{arXiv:1308.1804 [math.CT]}

\bibitem{Xu}
 J.~Xu.
   Flat covers of modules.
\textit{Lecture Notes in Math.}\ \textbf{1634}, Springer, 1996.

\bibitem{Zhu}
 X.~Zhu.
   Resolving resolution dimensions.
\textit{Algebras and Represent.\ Theory} \textbf{16}, \#4,
p.~1165--1191, 2013.

\end{thebibliography}

\end{document}
